% In this file you should put the actual content of the blueprint.
% It will be used both by the web and the print version.
% It should *not* include the \begin{document}
%
% If you want to split the blueprint content into several files then
% the current file can be a simple sequence of \input. Otherwise It
% can start with a \section or \chapter for instance.

% \input{theorems/gauss-lucas}

% \input{theorems/test}

\input{theorems/Abel's Theorem}

\input{theorems/Abel–Ruffini theorem}

\input{theorems/lovasz-local-lemma}

\input{theorems/brunn–minkowski theorem}

\input{theorems/beatty-sequence}

\input{theorems/Jacobson-Morozov Theorem}

\chapter{Bauer–Fike Theorem}

\begin{definition}[Condition Number of a Matrix]
    \label{def:Condition_Number}
    Let $X \in \mathbb{C}^{n \times n}$ be an invertible matrix. Its condition number in the $p$-norm, denoted by $\kappa_p(X)$, is defined as:
    \[ \kappa_p(X) = \|X\|_p \|X^{-1}\|_p \]
    where $\|\cdot\|_p$ is an induced matrix norm.
\end{definition}

\begin{lemma}[P-Norm of a Diagonal Matrix]
    \label{lem:P_Norm_of_Diagonal_Matrix}
    Let $D = \text{diag}(d_1, \dots, d_n) \in \mathbb{C}^{n \times n}$ be a diagonal matrix. For any induced matrix $p$-norm, where $1 \le p \le \infty$, its value is given by $\|D\|_p = \max_{1 \le i \le n} |d_i|$.
\end{lemma}

\begin{proof}
    By the definition of an induced matrix norm, $\|D\|_p = \sup_{x \neq 0} \frac{\|Dx\|_p}{\|x\|_p}$.
    For any vector $x = (x_1, \dots, x_n)^T$, the vector $Dx$ is $(d_1x_1, \dots, d_nx_n)^T$.
    The $p$-norm of $Dx$ is $\|Dx\|_p = \left( \sum_{i=1}^n |d_i x_i|^p \right)^{1/p}$.
    We can bound this as follows:
    \[ \sum_{i=1}^n |d_i x_i|^p \le \sum_{i=1}^n (\max_{j} |d_j|)^p |x_i|^p = (\max_{j} |d_j|)^p \sum_{i=1}^n |x_i|^p = (\max_{j} |d_j|)^p \|x\|_p^p \]
    Taking the $p$-th root, we get $\|Dx\|_p \le (\max_{j} |d_j|) \|x\|_p$. Thus, $\|D\|_p \le \max_{j} |d_j|$.
    To show equality, let $k$ be an index such that $|d_k| = \max_{j} |d_j|$. Consider the standard basis vector $x = e_k$, for which $\|e_k\|_p = 1$.
    Then $De_k$ is a vector with $d_k$ in the $k$-th position and zeros elsewhere, so $\|De_k\|_p = |d_k|$.
    This shows that the supremum is at least $|d_k|$, so $\|D\|_p \ge |d_k| = \max_{j} |d_j|$.
    Combining the inequalities, we conclude that $\|D\|_p = \max_{i} |d_i|$.
\end{proof}

\begin{lemma}[Eigenvalue Bound by Matrix Norm]
    \label{lem:Eigenvalue_Bound_by_Norm}
    Let $M \in \mathbb{C}^{n \times n}$ and let $\|\cdot\|$ be any matrix norm induced by a vector norm. If $\eta$ is an eigenvalue of $M$, then its modulus is bounded by the norm of the matrix: $|\eta| \le \|M\|$.
\end{lemma}

\begin{proof}
    Let $\eta$ be an eigenvalue of $M$ with a corresponding eigenvector $v \neq 0$. By definition, $Mv = \eta v$.
    Let $\|\cdot\|$ also denote the vector norm that induces the matrix norm. Taking the norm of both sides gives $\|Mv\| = \|\eta v\|$.
    Using the properties of norms, we have $\|\eta v\| = |\eta| \|v\|$. By the definition of an induced matrix norm, $\|Mv\| \le \|M\| \|v\|$.
    Combining these, we get $|\eta| \|v\| \le \|M\| \|v\|$. Since $v$ is an eigenvector, it is non-zero, so $\|v\| > 0$. We can divide by $\|v\|$ to obtain the result: $|\eta| \le \|M\|$.
\end{proof}

\begin{lemma}[Sub-multiplicativity of Induced Matrix Norms]
    \label{lem:Sub_Multiplicativity_of_Matrix_Norms}
    For any induced matrix norm $\|\cdot\|_p$, and for any matrices $A, B \in \mathbb{C}^{n \times n}$, the following inequality holds: $\|AB\|_p \le \|A\|_p \|B\|_p$.
\end{lemma}

\begin{proof}
    By the definition of an induced norm, for any vector $x \in \mathbb{C}^n$, we have $\|Ax\|_p \le \|A\|_p \|x\|_p$.
    Applying this property twice:
    \[ \|ABx\|_p = \|A(Bx)\|_p \le \|A\|_p \|Bx\|_p \le \|A\|_p (\|B\|_p \|x\|_p) = (\|A\|_p \|B\|_p) \|x\|_p \]
    By the definition of the induced norm for $AB$:
    \[ \|AB\|_p = \sup_{x \neq 0} \frac{\|ABx\|_p}{\|x\|_p} \le \sup_{x \neq 0} \frac{(\|A\|_p \|B\|_p) \|x\|_p}{\|x\|_p} = \|A\|_p \|B\|_p \]
    Thus, $\|AB\|_p \le \|A\|_p \|B\|_p$.
\end{proof}

\begin{lemma}[Determinant of a Matrix Product]
    \label{lem:Determinant_of_Product}
    For any square matrices $A, B \in \mathbb{C}^{n \times n}$, the determinant of their product is the product of their determinants: $\det(AB) = \det(A)\det(B)$. Consequently, a product of matrices is singular if and only if at least one of the matrices is singular.
\end{lemma}

\begin{proof}
    This is a standard result from linear algebra. If $A$ is singular, then $\det(A) = 0$. The rank of $AB$ is no greater than the rank of $A$, so $AB$ is also singular, meaning $\det(AB) = 0$. The equality holds.
    If $A$ is invertible, it can be written as a product of elementary matrices $A = E_k \dots E_1$. For any elementary matrix $E$ and any matrix $B$, $\det(EB) = \det(E)\det(B)$. Applying this repeatedly gives $\det(AB) = \det(E_k \dots E_1 B) = \det(E_k) \dots \det(E_1) \det(B) = \det(A) \det(B)$.
    Therefore, if a product $M_1 \dots M_k$ is singular, its determinant is 0, which implies $\det(M_1)\dots\det(M_k)=0$. This occurs if and only if $\det(M_i)=0$ for at least one $i$.
\end{proof}

\begin{lemma}[Similarity Invariance of the Determinant]
    \label{lem:Similarity_Invariance_of_Determinant}
    \uses{lem:Determinant_of_Product}
    The determinant is invariant under similarity transformations. For any matrix $M \in \mathbb{C}^{n \times n}$ and any invertible matrix $S \in \mathbb{C}^{n \times n}$, we have $\det(S^{-1}MS) = \det(M)$.
\end{lemma}

\begin{proof}
    \uses{lem:Determinant_of_Product}
    Using the multiplicative property of determinants from lemma \ref{lem:Determinant_of_Product}:
    \[ \det(S^{-1}MS) = \det(S^{-1}) \det(M) \det(S) \]
    Since $S$ is invertible, $SS^{-1}=I$, so $\det(S)\det(S^{-1})=\det(I)=1$. This implies $\det(S^{-1}) = \frac{1}{\det(S)}$.
    Substituting this back gives:
    \[ \det(S^{-1}MS) = \left(\frac{1}{\det(S)}\right) \det(M) \det(S) = \det(M) \]
\end{proof}

\begin{lemma}[Eigenvalues of a Diagonalizable Matrix and Invertibility]
    \label{lem:Eigenvalues_and_Invertibility}
    Let $A \in \mathbb{C}^{n \times n}$ be a diagonalizable matrix with eigendecomposition $A=V\Lambda V^{-1}$. The set of eigenvalues of $A$ is precisely the set of diagonal entries of $\Lambda$. Consequently, for any scalar $\mu \in \mathbb{C}$, the matrix $A - \mu I$ is invertible if and only if $\mu$ is not an eigenvalue of $A$.
\end{lemma}

\begin{proof}
    The matrix $A$ is similar to $\Lambda$, so they have the same eigenvalues. The eigenvalues of a diagonal matrix $\Lambda$ are its diagonal entries.
    A matrix $M$ is invertible if and only if $\det(M) \neq 0$. The condition $\det(A - \mu I) = 0$ is the characteristic equation for the eigenvalues of $A$. Thus, $A - \mu I$ is singular (not invertible) if and only if $\mu$ is an eigenvalue of $A$.
\end{proof}

\begin{lemma}[Eigenvalue Perturbation Identity]
    \label{lem:Eigenvalue_Perturbation_Identity}
    \uses{lem:Determinant_of_Product}
    \uses{lem:Similarity_Invariance_of_Determinant}
    \uses{lem:Eigenvalues_and_Invertibility}
    Let $A, \delta A \in \mathbb{C}^{n \times n}$, where $A$ is a diagonalizable matrix with $A = V\Lambda V^{-1}$. Let $\mu$ be an eigenvalue of $A + \delta A$ such that $\mu$ is not an eigenvalue of $A$. Then $-1$ is an eigenvalue of the matrix $(\Lambda - \mu I)^{-1} V^{-1} \delta A V$.
\end{lemma}

\begin{proof}
    \uses{lem:Determinant_of_Product}
    \uses{lem:Similarity_Invariance_of_Determinant}
    \uses{lem:Eigenvalues_and_Invertibility}
    Since $\mu$ is an eigenvalue of $A + \delta A$, we have $\det(A + \delta A - \mu I) = 0$. Using the fact that $A = V\Lambda V^{-1}$ and applying lemma \ref{lem:Similarity_Invariance_of_Determinant}, we can write:
    \begin{align*}
        0 &= \det(A + \delta A - \mu I) \\
          &= \det\left(V^{-1}(A + \delta A - \mu I)V\right) \\
          &= \det\left(V^{-1}AV + V^{-1}\delta A V - \mu I\right) \\
          &= \det\left(\Lambda + V^{-1}\delta A V - \mu I\right) \\
          &= \det\left((\Lambda - \mu I) + V^{-1}\delta A V\right)
    \end{align*}
    From lemma \ref{lem:Eigenvalues_and_Invertibility}, since $\mu$ is not an eigenvalue of $A$, the matrix $\Lambda - \mu I$ is invertible. We can factor it out:
    \[ \det\left((\Lambda - \mu I) \left(I + (\Lambda - \mu I)^{-1}V^{-1}\delta A V\right)\right) = 0 \]
    Using lemma \ref{lem:Determinant_of_Product}, this becomes:
    \[ \det(\Lambda - \mu I) \det\left(I + (\Lambda - \mu I)^{-1}V^{-1}\delta A V\right) = 0 \]
    Since $\Lambda - \mu I$ is invertible, $\det(\Lambda - \mu I) \neq 0$. Therefore, we must have:
    \[ \det\left(I + (\Lambda - \mu I)^{-1}V^{-1}\delta A V\right) = 0 \]
    This is the characteristic equation for the matrix $(\Lambda - \mu I)^{-1}V^{-1}\delta A V$ with eigenvalue $-1$.
\end{proof}

\begin{lemma}[Norm of Inverse Diagonal Matrix]
    \label{lem:Norm_of_Inverse_Diagonal_Matrix}
    \uses{lem:P_Norm_of_Diagonal_Matrix}
    Let $\Lambda = \text{diag}(\lambda_1, \dots, \lambda_n)$ be a diagonal matrix and $\mu \in \mathbb{C}$ such that $\mu \neq \lambda_i$ for all $i=1, \dots, n$. Then for any $p$-norm, the norm of the matrix $(\Lambda - \mu I)^{-1}$ is given by:
    \[ \|(\Lambda - \mu I)^{-1}\|_p = \frac{1}{\min_{\lambda \in \Lambda(A)} |\lambda - \mu|} \]
\end{lemma}

\begin{proof}
    \uses{lem:P_Norm_of_Diagonal_Matrix}
    The matrix $(\Lambda - \mu I)^{-1}$ is a diagonal matrix with diagonal entries $\frac{1}{\lambda_i - \mu}$. Applying lemma \ref{lem:P_Norm_of_Diagonal_Matrix}, its $p$-norm is the maximum of the absolute values of its diagonal entries:
    \[ \|(\Lambda - \mu I)^{-1}\|_p = \max_{i} \left| \frac{1}{\lambda_i - \mu} \right| = \frac{1}{\min_{i} |\lambda_i - \mu|} = \frac{1}{\min_{\lambda \in \Lambda(A)} |\lambda - \mu|} \]
\end{proof}

\begin{theorem}[Bauer–Fike Theorem]
    \label{thm:Bauer-Fike}
    \uses{def:Condition_Number}
    Let $A \in \mathbb{C}^{n \times n}$ be a diagonalizable matrix, and let $\mu$ be an eigenvalue of the perturbed matrix $A + \delta A$. Then there exists an eigenvalue $\lambda$ of $A$ such that:
    \[ |\lambda - \mu| \leq \kappa_p(V) \|\delta A\|_p \]
    where $V$ is the matrix of eigenvectors of $A$.
\end{theorem}

\begin{proof}
    \uses{def:Condition_Number}
    \uses{lem:Eigenvalue_Bound_by_Norm}
    \uses{lem:Sub_Multiplicativity_of_Matrix_Norms}
    \uses{lem:Eigenvalue_Perturbation_Identity}
    \uses{lem:Norm_of_Inverse_Diagonal_Matrix}
    If $\mu$ is an eigenvalue of $A$, the inequality holds trivially as $|\mu - \mu| = 0$.
    Assume $\mu$ is not an eigenvalue of $A$. By lemma \ref{lem:Eigenvalue_Perturbation_Identity}, $-1$ is an eigenvalue of $M = (\Lambda - \mu I)^{-1} V^{-1} \delta A V$.
    By lemma \ref{lem:Eigenvalue_Bound_by_Norm}, we have $1 = |-1| \leq \|M\|_p$.
    Using the sub-multiplicative property from lemma \ref{lem:Sub_Multiplicativity_of_Matrix_Norms}:
    \[ 1 \leq \|(\Lambda - \mu I)^{-1}\|_p \|V^{-1}\|_p \|\delta A\|_p \|V\|_p \]
    Using definition \ref{def:Condition_Number}, this becomes $1 \leq \|(\Lambda - \mu I)^{-1}\|_p \kappa_p(V) \|\delta A\|_p$.
    Applying lemma \ref{lem:Norm_of_Inverse_Diagonal_Matrix}:
    \[ 1 \leq \frac{1}{\min_{\lambda \in \Lambda(A)} |\lambda - \mu|} \kappa_p(V) \|\delta A\|_p \]
    Rearranging gives the desired result: $\min_{\lambda \in \Lambda(A)} |\lambda - \mu| \leq \kappa_p(V) \|\delta A\|_p$.
\end{proof}

\chapter{Casorati–Weierstrass Theorem}

\begin{definition}[Essential Singularity]
    \label{def:Essential_Singularity}
    Let $f$ be a function that is holomorphic on a punctured neighborhood $D = U \setminus \{z_0\}$ of a point $z_0$. The point $z_0$ is called an essential singularity of $f$ if it is neither a removable singularity nor a pole.
\end{definition}

\begin{definition}[Dense Set]
    \label{def:Dense_Set}
    A subset $A$ of the complex plane $\mathbb{C}$ is said to be dense in $\mathbb{C}$ if for every point $w \in \mathbb{C}$ and for every $\varepsilon > 0$, there exists a point $a \in A$ such that $|a - w| < \varepsilon$. Equivalently, the closure of $A$ is $\mathbb{C}$.
\end{definition}

\begin{lemma}[Analyticity of Holomorphic Functions]
    \label{lem:Holomorphic_Implies_Analytic}
    If a function $g$ is holomorphic in an open disk $D(z_0, R)$, then it is analytic at $z_0$. This means $g$ can be represented by its Taylor series expansion, $g(z) = \sum_{n=0}^{\infty} a_n (z-z_0)^n$, for all $z \in D(z_0, R)$.
\end{lemma}

\begin{proof}
    This is a cornerstone theorem of complex analysis, following from Cauchy's integral formula for derivatives. The proof is omitted.
\end{proof}

\begin{lemma}[Riemann's Theorem on Removable Singularities]
    \label{lem:Riemann_Removable_Singularity}
    Let $f$ be a function that is holomorphic and bounded in a punctured neighborhood $U \setminus \{z_0\}$. Then the singularity at $z_0$ is removable, meaning $f$ can be extended to a holomorphic function on all of $U$.
\end{lemma}

\begin{proof}
    This is a fundamental result in complex analysis and its proof is omitted here.
\end{proof}

\begin{lemma}[Characterization of Removable Singularities by Limit]
    \label{lem:Characterization_of_Removable_Singularities}
    Let $f$ be holomorphic on a punctured neighborhood of $z_0$. The singularity at $z_0$ is removable if and only if the limit $\lim_{z \to z_0} f(z)$ exists and is finite.
\end{lemma}

\begin{proof}
    This is a standard characterization of removable singularities, presented here without proof.
\end{proof}

\begin{lemma}[Characterization of Poles by Limit]
    \label{lem:Characterization_of_Poles}
    Let $f$ be holomorphic on a punctured neighborhood of $z_0$. The singularity at $z_0$ is a pole if and only if $\lim_{z \to z_0} |f(z)| = \infty$.
\end{lemma}

\begin{proof}
    This is a standard characterization of poles, presented here without proof.
\end{proof}

\begin{lemma}[Algebra of Limits for Complex Functions]
    \label{lem:Algebra_of_Limits}
    Let $g(z)$ be a complex function defined on a punctured neighborhood of $z_0$. If $\lim_{z \to z_0} g(z) = L$, where $L \in \mathbb{C}$ and $L \neq 0$, then for any constant $b \in \mathbb{C}$, the limit of $\frac{1}{g(z)} + b$ as $z \to z_0$ exists and is equal to $\frac{1}{L} + b$.
\end{lemma}

\begin{proof}
    This follows from the standard properties of limits concerning sums, constants, and reciprocals, which hold for complex-valued functions. The proof is omitted.
\end{proof}

\begin{lemma}[Limit of Reciprocal Function]
    \label{lem:Limit_of_Reciprocal}
    Let $d(z)$ be a complex function defined on a punctured neighborhood of $z_0$ such that $\lim_{z \to z_0} d(z) = 0$. Then $\lim_{z \to z_0} \frac{1}{|d(z)|} = \infty$, and consequently, for any constant $b \in \mathbb{C}$, $\lim_{z \to z_0} \left|\frac{1}{d(z)} + b\right| = \infty$.
\end{lemma}

\begin{proof}
    This is a standard property of limits involving infinity. As $z \to z_0$, $|d(z)|$ becomes arbitrarily small and positive, so its reciprocal becomes arbitrarily large. Adding a constant $b$ does not change this behavior for the magnitude of the limit.
\end{proof}

\begin{lemma}[Factorization at a Zero]
    \label{lem:Zero_of_Holomorphic_Function}
    \uses{lem:Holomorphic_Implies_Analytic}
    Let $g$ be a holomorphic function in a neighborhood of $z_0$ such that $g(z_0) = 0$. If $g$ is not identically zero in this neighborhood, then there exists a unique integer $m \geq 1$ (the order of the zero) and a holomorphic function $h$ defined in a neighborhood of $z_0$ such that $h(z_0) \neq 0$ and $g(z) = (z-z_0)^m h(z)$.
\end{lemma}

\begin{proof}
    \uses{lem:Holomorphic_Implies_Analytic}
    By lemma \ref{lem:Holomorphic_Implies_Analytic}, $g$ has a Taylor series expansion $g(z) = \sum_{n=0}^{\infty} a_n (z-z_0)^n$. Since $g(z_0)=0$, $a_0=0$. As $g$ is not identically zero, not all coefficients are zero. Let $m \geq 1$ be the index of the first non-zero coefficient, so $a_m \neq 0$ and $a_k=0$ for $k < m$. We can then write $g(z) = (z-z_0)^m \sum_{n=m}^{\infty} a_n (z-z_0)^{n-m} = (z-z_0)^m h(z)$, where $h(z) = \sum_{k=0}^{\infty} a_{m+k} (z-z_0)^k$. The function $h(z)$ is holomorphic and $h(z_0) = a_m \neq 0$.
\end{proof}

\begin{lemma}[Existence of a Bounded Auxiliary Function]
    \label{lem:Bounded_Auxiliary_Function}
    Let $f$ be a function that is holomorphic on a punctured neighborhood $V \setminus \{z_0\}$. If the image $f(V \setminus \{z_0\})$ is not dense in $\mathbb{C}$, then there exists a function $g(z)$ which is holomorphic and bounded on $V \setminus \{z_0\}$.
\end{lemma}

\begin{proof}
    If the image $f(V \setminus \{z_0\})$ is not dense in $\mathbb{C}$, then there exists a point $b \in \mathbb{C}$ and an open disk $D(b, \varepsilon)$ with radius $\varepsilon > 0$ that is disjoint from the image. This means $|f(z) - b| \geq \varepsilon$ for all $z \in V \setminus \{z_0\}$.
    Define $g(z) = \frac{1}{f(z) - b}$. Since $f(z) - b$ is holomorphic and never zero on $V \setminus \{z_0\}$, $g(z)$ is also holomorphic there. Furthermore, $|g(z)| = \frac{1}{|f(z) - b|} \leq \frac{1}{\varepsilon}$, so $g(z)$ is bounded.
\end{proof}

\begin{lemma}[Classification of the Singularity of f]
    \label{lem:Classification_of_f}
    \uses{lem:Bounded_Auxiliary_Function}
    \uses{lem:Riemann_Removable_Singularity}
    \uses{lem:Characterization_of_Removable_Singularities}
    \uses{lem:Characterization_of_Poles}
    \uses{lem:Algebra_of_Limits}
    \uses{lem:Limit_of_Reciprocal}
    \uses{lem:Zero_of_Holomorphic_Function}
    If the image $f(V \setminus \{z_0\})$ is not dense in $\mathbb{C}$, then the singularity of $f$ at $z_0$ is either removable or a pole.
\end{lemma}

\begin{proof}
    \uses{lem:Bounded_Auxiliary_Function}
    \uses{lem:Riemann_Removable_Singularity}
    \uses{lem:Characterization_of_Removable_Singularities}
    \uses{lem:Characterization_of_Poles}
    \uses{lem:Algebra_of_Limits}
    \uses{lem:Limit_of_Reciprocal}
    \uses{lem:Zero_of_Holomorphic_Function}
    By lemma \ref{lem:Bounded_Auxiliary_Function}, we can construct a bounded, holomorphic function $g(z) = \frac{1}{f(z) - b}$ on $V \setminus \{z_0\}$. By lemma \ref{lem:Riemann_Removable_Singularity}, $g(z)$ has a removable singularity at $z_0$, so it can be extended to a holomorphic function on $V$. Let $L = g(z_0) = \lim_{z \to z_0} g(z)$.

    Case 1: $L \neq 0$. Using lemma \ref{lem:Algebra_of_Limits} on $f(z) = \frac{1}{g(z)} + b$, we find $\lim_{z \to z_0} f(z) = \frac{1}{L} + b$, which is finite. By lemma \ref{lem:Characterization_of_Removable_Singularities}, $f$ has a removable singularity at $z_0$.

    Case 2: $L = 0$. Since $g$ is not identically zero, by lemma \ref{lem:Zero_of_Holomorphic_Function}, we can write $g(z) = (z-z_0)^m h(z)$ where $m \geq 1$ and $h(z_0) \neq 0$. Then $f(z) = \frac{1}{(z-z_0)^m h(z)} + b$. The term $(z-z_0)^m h(z)$ approaches $0$ as $z \to z_0$. By lemma \ref{lem:Limit_of_Reciprocal}, $\lim_{z \to z_0} |f(z)| = \infty$. By lemma \ref{lem:Characterization_of_Poles}, $f$ has a pole at $z_0$.

    In either case, the singularity of $f$ is not essential.
\end{proof}

\begin{theorem}[Casorati–Weierstrass Theorem]
    \label{thm:Casorati-Weierstrass_Theorem}
    \uses{def:Essential_Singularity}
    \uses{def:Dense_Set}
    Let $f$ be a function that is holomorphic on a punctured neighborhood $V \setminus \{z_0\}$ and has an essential singularity at $z_0$. Then the image of $V \setminus \{z_0\}$ under $f$, denoted $f(V \setminus \{z_0\})$, is dense in the complex plane $\mathbb{C}$.
\end{theorem}

\begin{proof}
    \uses{def:Essential_Singularity}
    \uses{lem:Classification_of_f}
    We argue by contradiction. Assume that the image $f(V \setminus \{z_0\})$ is not dense in $\mathbb{C}$.
    
    According to lemma \ref{lem:Classification_of_f}, if the image is not dense, then the singularity of $f$ at $z_0$ must be either a removable singularity or a pole.
    
    This conclusion contradicts the initial hypothesis that $z_0$ is an essential singularity of $f$ (see definition \ref{def:Essential_Singularity}). Therefore, the assumption must be false, which means the image $f(V \setminus \{z_0\})$ must be dense in $\mathbb{C}$.
\end{proof}

\chapter{Ceva's Theorem}

\begin{definition}[Cevian]
    \label{def:Cevian}
    Given a triangle $\triangle ABC$, a cevian is a line segment that joins a vertex to a point on the opposite side. Let the lines $AO$, $BO$, $CO$ be drawn from the vertices to a common point $O$ (not on one of the sides of $\triangle ABC$), to meet opposite sides at $D$, $E$, $F$ respectively. The segments $AD$, $BE$, and $CF$ are known as the cevians of the point $O$.
\end{definition}

\begin{lemma}[Property of Proportions]
    \label{lem:Property_of_Proportions}
    Let $a, b, c, d$ be real numbers with $b \neq 0$, $d \neq 0$ and $b \neq d$. If $\frac{a}{b} = \frac{c}{d} = k$, then $\frac{a-c}{b-d} = k$.
\end{lemma}

\begin{proof}
    From the given condition $\frac{a}{b} = k$ and $\frac{c}{d} = k$, we can write $a = kb$ and $c = kd$.
    Substituting these into the expression $\frac{a-c}{b-d}$ gives:
    \[ \frac{a-c}{b-d} = \frac{kb - kd}{b-d} = \frac{k(b-d)}{b-d} \]
    Since $b \neq d$, the term $(b-d)$ is non-zero, so we can cancel it from the numerator and denominator, which yields:
    \[ \frac{k(b-d)}{b-d} = k \]
    Thus, $\frac{a-c}{b-d} = k = \frac{a}{b} = \frac{c}{d}$.
\end{proof}

\begin{lemma}[Area Addition Postulate]
    \label{lem:Area_Addition_Postulate}
    Let $A$ and $D$ be two distinct points, and let $O$ be a point in the interior of the line segment $AD$. For any point $B$ not on the line $AD$, the area of $\triangle ABD$ is the sum of the areas of $\triangle ABO$ and $\triangle OBD$. That is, $|\triangle ABD| = |\triangle ABO| + |\triangle OBD|$.
\end{lemma}

\begin{proof}
    The line segment $BO$ partitions the triangle $\triangle ABD$ into two smaller, non-overlapping triangles: $\triangle ABO$ and $\triangle OBD$. The base $AD$ of $\triangle ABD$ is the sum of the bases $AO$ and $OD$ of $\triangle ABO$ and $\triangle OBD$ respectively. Let $h_B$ be the perpendicular distance from vertex $B$ to the line $AD$. The areas are:
    $|\triangle ABD| = \frac{1}{2} |\overline{AD}| \cdot h_B$
    $|\triangle ABO| = \frac{1}{2} |\overline{AO}| \cdot h_B$
    $|\triangle OBD| = \frac{1}{2} |\overline{OD}| \cdot h_B$
    Since $|\overline{AD}| = |\overline{AO}| + |\overline{OD}|$, it follows that
    \[ |\triangle ABD| = \frac{1}{2} (|\overline{AO}| + |\overline{OD}|) \cdot h_B = \frac{1}{2} |\overline{AO}| \cdot h_B + \frac{1}{2} |\overline{OD}| \cdot h_B = |\triangle ABO| + |\triangle OBD| \]
\end{proof}

\begin{lemma}[Area Subtraction Property]
    \label{lem:Area_Subtraction_Property}
    \uses{def:Cevian}
    \uses{lem:Area_Addition_Postulate}
    For a triangle $\triangle ABC$, let $D$ be a point on the line containing side $BC$, and let $O$ be a point on the cevian $AD$. The area of $\triangle ABO$ is the difference between the areas of $\triangle ABD$ and $\triangle OBD$. That is, $|\triangle ABO| = |\triangle ABD| - |\triangle OBD|$. Similarly, $|\triangle CAO| = |\triangle ACD| - |\triangle OCD|$.
\end{lemma}

\begin{proof}
    \uses{lem:Area_Addition_Postulate}
    By lemma \ref{lem:Area_Addition_Postulate}, since point $O$ lies on the line segment $AD$, the area of the larger triangle $\triangle ABD$ is the sum of the areas of the two triangles it is partitioned into:
    \[ |\triangle ABD| = |\triangle ABO| + |\triangle OBD| \]
    Rearranging the terms, we get:
    \[ |\triangle ABO| = |\triangle ABD| - |\triangle OBD| \]
    Similarly, for $\triangle ACD$:
    \[ |\triangle ACD| = |\triangle CAO| + |\triangle OCD| \]
    Rearranging gives:
    \[ |\triangle CAO| = |\triangle ACD| - |\triangle OCD| \]
\end{proof}

\begin{lemma}[Ratio of Areas of Triangles with a Common Altitude]
    \label{lem:Ratio_of_Areas_of_Triangles_with_a_Common_Altitude}
    For a triangle $\triangle ABC$ and a point $D$ on the line containing $BC$, the ratio of the signed lengths of the segments $BD$ and $DC$ is equal to the ratio of the areas of the triangles $\triangle ABD$ and $\triangle ACD$. That is,
    \[ \frac{\overline{BD}}{\overline{DC}} = \frac{|\triangle ABD|}{|\triangle ACD|} \]
    Similarly, for any point $O$ not on the line $BC$,
    \[ \frac{\overline{BD}}{\overline{DC}} = \frac{|\triangle OBD|}{|\triangle OCD|} \]
\end{lemma}

\begin{proof}
    The area of a triangle is given by $\frac{1}{2} \times \text{base} \times \text{height}$. Let $h$ be the length of the altitude from vertex $A$ to the line containing side $BC$. Then the areas of $\triangle ABD$ and $\triangle ACD$ are $|\triangle ABD| = \frac{1}{2} |\overline{BD}| \cdot h$ and $|\triangle ACD| = \frac{1}{2} |\overline{DC}| \cdot h$.
    The ratio of their areas, considering signed lengths for the bases, is
    \[ \frac{|\triangle ABD|}{|\triangle ACD|} = \frac{\frac{1}{2} \overline{BD} \cdot h}{\frac{1}{2} \overline{DC} \cdot h} = \frac{\overline{BD}}{\overline{DC}} \]
    A similar argument holds for the triangles $\triangle OBD$ and $\triangle OCD$ with a common vertex $O$ and bases on the same line $BC$. The altitude from $O$ to line $BC$ is common for both triangles.
\end{proof}

\begin{lemma}[Segment Ratio as a Ratio of Triangle Areas]
    \label{lem:Segment_Ratio_as_a_Ratio_of_Triangle_Areas}
    \uses{def:Cevian}
    \uses{lem:Property_of_Proportions}
    \uses{lem:Area_Subtraction_Property}
    \uses{lem:Ratio_of_Areas_of_Triangles_with_a_Common_Altitude}
    In $\triangle ABC$, let $AD$ be a cevian from $A$ to $BC$, and let $O$ be a point on $AD$. The ratio of the segments $\overline{BD}$ and $\overline{DC}$ is given by the ratio of the areas of $\triangle ABO$ and $\triangle CAO$.
    \[ \frac{\overline{BD}}{\overline{DC}} = \frac{|\triangle ABO|}{|\triangle CAO|} \]
\end{lemma}

\begin{proof}
    \uses{lem:Ratio_of_Areas_of_Triangles_with_a_Common_Altitude}
    \uses{lem:Property_of_Proportions}
    \uses{lem:Area_Subtraction_Property}
    From lemma \ref{lem:Ratio_of_Areas_of_Triangles_with_a_Common_Altitude}, we have two expressions for the ratio $\frac{\overline{BD}}{\overline{DC}}$:
    \[ \frac{\overline{BD}}{\overline{DC}} = \frac{|\triangle ABD|}{|\triangle ACD|} \quad \text{and} \quad \frac{\overline{BD}}{\overline{DC}} = \frac{|\triangle OBD|}{|\triangle OCD|} \]
    This implies $\frac{|\triangle ABD|}{|\triangle ACD|} = \frac{|\triangle OBD|}{|\triangle OCD|}$.
    By lemma \ref{lem:Property_of_Proportions}, we can write:
    \[ \frac{|\triangle ABD|}{|\triangle ACD|} = \frac{|\triangle ABD| - |\triangle OBD|}{|\triangle ACD| - |\triangle OCD|} \]
    Using lemma \ref{lem:Area_Subtraction_Property}, we can substitute the numerators and denominators:
    \[ |\triangle ABD| - |\triangle OBD| = |\triangle ABO| \]
    \[ |\triangle ACD| - |\triangle OCD| = |\triangle CAO| \]
    Therefore,
    \[ \frac{\overline{BD}}{\overline{DC}} = \frac{|\triangle ABO|}{|\triangle CAO|} \]
\end{proof}

\begin{theorem}[Ceva's Theorem (Direct)]
    \label{thm:Cevas_Theorem_Direct}
    \uses{def:Cevian}
    \uses{lem:Segment_Ratio_as_a_Ratio_of_Triangle_Areas}
    Let $\triangle ABC$ be a triangle, and let $D, E, F$ be points on the lines containing sides $BC, CA, AB$ respectively. If the cevians $AD, BE, CF$ are concurrent at a point $O$, then
    \[ \frac{\overline{AF}}{\overline{FB}} \cdot \frac{\overline{BD}}{\overline{DC}} \cdot \frac{\overline{CE}}{\overline{EA}} = 1 \]
\end{theorem}

\begin{proof}
    \uses{lem:Segment_Ratio_as_a_Ratio_of_Triangle_Areas}
    Since the cevians are concurrent at $O$, we can apply lemma \ref{lem:Segment_Ratio_as_a_Ratio_of_Triangle_Areas} to each cevian:
    \[ \frac{\overline{BD}}{\overline{DC}} = \frac{|\triangle ABO|}{|\triangle CAO|} \]
    \[ \frac{\overline{CE}}{\overline{EA}} = \frac{|\triangle BCO|}{|\triangle ABO|} \]
    \[ \frac{\overline{AF}}{\overline{FB}} = \frac{|\triangle CAO|}{|\triangle BCO|} \]
    Multiplying these three equations together yields:
    \[ \frac{\overline{AF}}{\overline{FB}} \cdot \frac{\overline{BD}}{\overline{DC}} \cdot \frac{\overline{CE}}{\overline{EA}} = \frac{|\triangle CAO|}{|\triangle BCO|} \cdot \frac{|\triangle ABO|}{|\triangle CAO|} \cdot \frac{|\triangle BCO|}{|\triangle ABO|} \]
    The terms in the numerator cancel with the terms in the denominator, resulting in:
    \[ \frac{\overline{AF}}{\overline{FB}} \cdot \frac{\overline{BD}}{\overline{DC}} \cdot \frac{\overline{CE}}{\overline{EA}} = 1 \]
\end{proof}

\begin{lemma}[Uniqueness of a Point Dividing a Segment]
    \label{lem:Uniqueness_of_Point_Dividing_a_Segment}
    Let $A$ and $B$ be two distinct points on a line $L$. For any real number $r$ such that $r \neq -1$, there exists a unique point $F$ on $L$ such that the ratio of signed lengths $\frac{\overline{AF}}{\overline{FB}} = r$.
\end{lemma}

\begin{proof}
    Let the line $L$ be the real number line, with $A$ at the origin (0) and $B$ at coordinate $b \neq 0$. Let $F$ be a point on $L$ with coordinate $x$. The signed lengths are $\overline{AF} = x - 0 = x$ and $\overline{FB} = b - x$. The condition is:
    \[ \frac{x}{b-x} = r \]
    We solve for $x$:
    \[ x = r(b-x) \implies x = rb - rx \implies x(1+r) = rb \]
    Since $r \neq -1$, we have $1+r \neq 0$, so we can divide by it:
    \[ x = \frac{rb}{1+r} \]
    This equation gives a single, unique value for the coordinate $x$. Therefore, the point $F$ is uniquely determined.
\end{proof}

\begin{lemma}[Ray Intersection Lemma]
    \label{lem:Ray_Intersection_Lemma}
    In a triangle $\triangle ABC$, if a line enters the triangle at vertex $C$ and passes through an interior point $O$, it must intersect the opposite side $AB$.
\end{lemma}

\begin{proof}
    This lemma is a consequence of Pasch's Axiom. The line passing through $C$ and $O$ divides the plane into two half-planes. The vertices $A$ and $B$ lie in opposite half-planes because the point $O$ is interior to $\triangle ABC$, meaning $O$ is on the same side of $AB$ as $C$, the same side of $BC$ as $A$, and the same side of $AC$ as $B$. The line $CO$ passes through vertex $C$ and interior point $O$. Since $A$ and $B$ are on opposite sides of the line $CO$, the segment $AB$ must intersect the line $CO$ at a point. This intersection point lies on the segment $AB$.
\end{proof}

\begin{theorem}[Ceva's Theorem (Converse)]
    \label{thm:Cevas_Theorem_Converse}
    \uses{thm:Cevas_Theorem_Direct}
    \uses{lem:Uniqueness_of_Point_Dividing_a_Segment}
    \uses{lem:Ray_Intersection_Lemma}
    Let $\triangle ABC$ be a triangle, and let $D, E, F$ be points on the lines containing sides $BC, CA, AB$ respectively. If
    \[ \frac{\overline{AF}}{\overline{FB}} \cdot \frac{\overline{BD}}{\overline{DC}} \cdot \frac{\overline{CE}}{\overline{EA}} = 1 \]
    then the lines $AD, BE, CF$ are concurrent.
\end{theorem}

\begin{proof}
    \uses{thm:Cevas_Theorem_Direct}
    \uses{lem:Uniqueness_of_Point_Dividing_a_Segment}
    \uses{lem:Ray_Intersection_Lemma}
    Let the lines $AD$ and $BE$ intersect at a point $O$. By lemma \ref{lem:Ray_Intersection_Lemma}, the line $CO$ must intersect the line $AB$ at some point. Let this point be $F'$.
    Since the cevians $AD$, $BE$, and $CF'$ are concurrent at $O$ by construction, we can apply the direct part of Ceva's theorem, theorem \ref{thm:Cevas_Theorem_Direct}, to these cevians:
    \[ \frac{\overline{AF'}}{\overline{F'B}} \cdot \frac{\overline{BD}}{\overline{DC}} \cdot \frac{\overline{CE}}{\overline{EA}} = 1 \]
    We are given the condition:
    \[ \frac{\overline{AF}}{\overline{FB}} \cdot \frac{\overline{BD}}{\overline{DC}} \cdot \frac{\overline{CE}}{\overline{EA}} = 1 \]
    Comparing the two equations, we can cancel the common non-zero terms to obtain:
    \[ \frac{\overline{AF'}}{\overline{F'B}} = \frac{\overline{AF}}{\overline{FB}} \]
    By lemma \ref{lem:Uniqueness_of_Point_Dividing_a_Segment}, there is only one point on the line $AB$ that divides the segment $AB$ in this specific ratio. Since both $F$ and $F'$ satisfy this condition, they must be the same point, i.e., $F=F'$.
    Therefore, the line $CF$ is the same as the line $CF'$, which passes through the point $O$. This means that the lines $AD, BE, CF$ are concurrent at $O$.
\end{proof}

\chapter{Chvátal's Art Gallery Theorem}

\begin{definition}[Guard Set]
    \label{def:Guard_Set}
    Let $P$ be a simple polygon. A set of points $S$ is said to be a guard set for $P$ if for every point $p \in P$, there exists a point $q \in S$ such that the line segment $\overline{pq}$ is entirely contained within $P$.
\end{definition}

\begin{definition}[Polygon Triangulation]
    \label{def:Polygon_Triangulation}
    A triangulation of a simple polygon $P$ with $n$ vertices is a decomposition of $P$ into $n-2$ triangles by a set of non-intersecting diagonals connecting pairs of vertices of $P$.
\end{definition}

\begin{definition}[3-Coloring of a Graph]
    \label{def:3-Coloring_of_a_Graph}
    A 3-coloring of a graph $G = (V, E)$ is a function $c: V \to \{1, 2, 3\}$ such that for any edge $(u, v) \in E$, we have $c(u) \neq c(v)$. A graph that admits a 3-coloring is called 3-colorable.
\end{definition}

\begin{lemma}[Finding a Specific Internal Diagonal]
    \label{lem:Finding_a_Specific_Internal_Diagonal}
    Let $v$ be a convex vertex of a simple polygon $P$ with adjacent vertices $u$ and $w$. If the interior of the triangle $\triangle uvw$ contains other vertices of $P$, let $p$ be one such vertex that is furthest from the line segment $\overline{uw}$. Then the line segment $\overline{vp}$ is an internal diagonal of $P$.
\end{lemma}

\begin{proof}
    The segment $\overline{vp}$ connects two vertices of the polygon. To be an internal diagonal, it must lie entirely in the interior of $P$, except for its endpoints. Assume for contradiction that $\overline{vp}$ intersects an edge of the polygon, say $\overline{ab}$. This edge $\overline{ab}$ cannot be incident to $v$. The line segment $\overline{vp}$ is contained within the triangle $\triangle uvw$. Therefore, the intersecting edge $\overline{ab}$ must have at least one endpoint, say $a$, inside $\triangle uvw$.
    
    The vertex $p$ was chosen as the vertex inside $\triangle uvw$ furthest from the line containing $\overline{uw}$. However, if we consider the triangle formed by $v$ and the line $\overline{uw}$, any vertex $a$ on an edge intersecting $\overline{vp}$ must lie further from $\overline{uw}$ than $p$ does. This contradicts the choice of $p$. Thus, no edge of the polygon can intersect the interior of $\overline{vp}$. Therefore, $\overline{vp}$ is an internal diagonal.
\end{proof}

\begin{lemma}[Existence of an Internal Diagonal]
    \label{lem:Existence_of_an_Internal_Diagonal}
    Any simple polygon with $n > 3$ vertices has at least one diagonal that lies entirely within its interior.
\end{lemma}

\begin{proof}
    \uses{lem:Finding_a_Specific_Internal_Diagonal}
    Let $v$ be a vertex of the polygon with a convex interior angle. Let $u$ and $w$ be the vertices adjacent to $v$. Consider the line segment $\overline{uw}$.
    
    Case 1: The segment $\overline{uw}$ lies entirely within the polygon and no other vertices lie on it. In this case, $\overline{uw}$ is an internal diagonal.
    
    Case 2: The segment $\overline{uw}$ does not lie entirely within the polygon, or other vertices lie on it. This implies that the interior of the triangle $\triangle uvw$ contains one or more vertices of the polygon. By lemma \ref{lem:Finding_a_Specific_Internal_Diagonal}, we can find a vertex $p$ inside $\triangle uvw$ such that $\overline{vp}$ is an internal diagonal.
    
    In either case, an internal diagonal exists. Since every simple polygon must have at least one convex vertex, this condition can always be met.
\end{proof}

\begin{lemma}[Polygon Splitting by Diagonal]
    \label{lem:Polygon_Splitting_by_Diagonal}
    Let $P$ be a simple polygon with $n$ vertices. An internal diagonal partitions $P$ into two smaller simple polygons, whose vertex counts, $n_1$ and $n_2$, satisfy $n_1 + n_2 = n+2$.
\end{lemma}
\begin{proof}
    Let the diagonal connect vertices $v_i$ and $v_j$. This diagonal splits the boundary of the original polygon into two polygonal chains. The first smaller polygon, $P_1$, is formed by the diagonal $\overline{v_i v_j}$ and the chain of vertices from $v_i$ to $v_j$. The second smaller polygon, $P_2$, is formed by the same diagonal and the remaining chain of vertices from $v_j$ back to $v_i$. The vertices $v_i$ and $v_j$ are part of both new polygons. All other $n-2$ vertices of the original polygon belong to exactly one of the new polygons. Therefore, the total number of vertices in the two new polygons is $(n-2) + 2 + 2 = n+2$.
\end{proof}

\begin{lemma}[Existence of Triangulation]
    \label{lem:Existence_of_Triangulation}
    \uses{def:Polygon_Triangulation}
    Any simple polygon with $n$ vertices can be triangulated.
\end{lemma}

\begin{proof}
    \uses{lem:Existence_of_an_Internal_Diagonal}
    \uses{lem:Polygon_Splitting_by_Diagonal}
    The proof is by induction on the number of vertices $n$. The base case is $n=3$, a triangle, which is already triangulated.
    For $n > 3$, lemma \ref{lem:Existence_of_an_Internal_Diagonal} guarantees the existence of an internal diagonal. This diagonal splits the polygon into two smaller simple polygons. By lemma \ref{lem:Polygon_Splitting_by_Diagonal}, these polygons have $n_1$ and $n_2$ vertices, where $n_1, n_2 < n$. By the induction hypothesis, both smaller polygons can be triangulated. The union of these two triangulations forms a triangulation of the original polygon.
\end{proof}

\begin{lemma}[Dual Graph of Triangulation is a Tree]
    \label{lem:Dual_Graph_is_a_Tree}
    \uses{def:Polygon_Triangulation}
    \uses{lem:Existence_of_Triangulation}
    Let $P$ be a simple polygon and let $T$ be a triangulation of $P$. The dual graph of $T$, which has a vertex for each triangle and an edge connecting vertices corresponding to adjacent triangles, is a tree.
\end{lemma}

\begin{proof}
    The dual graph is connected because the polygon is a single connected region. Any path between two triangles in the triangulation corresponds to a path in the dual graph. A cycle in the dual graph would correspond to a set of triangles that enclose a region not belonging to the polygon, forming a hole. Since $P$ is a simple polygon, it has no holes. Therefore, the dual graph is connected and acyclic, which means it is a tree.
\end{proof}

\begin{lemma}[Existence of Leaves in a Tree]
    \label{lem:Existence_of_Leaves_in_a_Tree}
    Any finite tree with at least two vertices has at least two leaves (vertices of degree 1).
\end{lemma}

\begin{proof}
    Consider a longest path in the tree. Let the endpoints of this path be $u$ and $v$. If the degree of $u$ were greater than 1, it would be adjacent to some vertex $w$ not on the path. The edge $(u,w)$ cannot create a cycle, as the graph is a tree. Thus, $w$ could be used to extend the path, contradicting the assumption that the path was the longest. Therefore, the degree of $u$ must be 1. The same argument applies to $v$.
\end{proof}

\begin{lemma}[Dual Graph Leaf corresponds to Polygon Ear]
    \label{lem:Dual_Graph_Leaf_is_Ear}
    \uses{def:Polygon_Triangulation}
    \uses{lem:Dual_Graph_is_a_Tree}
    Let $\mathcal{T}$ be a triangulation of a simple polygon $P$, and let $G$ be the dual graph of $\mathcal{T}$. A vertex of $G$ is a leaf if and only if the corresponding triangle in $\mathcal{T}$ is an "ear" of $P$ (a triangle formed by two polygon edges and one internal diagonal).
\end{lemma}

\begin{proof}
    A vertex in the dual graph $G$ corresponds to a triangle in the triangulation $\mathcal{T}$. The degree of a vertex in $G$ is the number of adjacent triangles, which is the number of sides of the corresponding triangle that are internal diagonals.
    
    If a vertex in $G$ is a leaf, its degree is 1. This means its corresponding triangle has exactly one side that is an internal diagonal, shared with another triangle. The other two sides cannot be diagonals, so they must be edges of the original polygon. This is the definition of a polygon ear.
    
    Conversely, if a triangle is an ear, it is bounded by two polygon edges and one internal diagonal. It shares only this one diagonal with the rest of the triangulation. Therefore, its corresponding vertex in the dual graph $G$ has degree 1 and is a leaf.
\end{proof}

\begin{lemma}[3-Colorability of a Triangulated Polygon]
    \label{lem:3-Colorability_of_a_Triangulated_Polygon}
    \uses{def:Polygon_Triangulation}
    \uses{def:3-Coloring_of_a_Graph}
    \uses{lem:Existence_of_Triangulation}
    The vertices of any triangulated simple polygon can be 3-colored.
\end{lemma}

\begin{proof}
    \uses{lem:Dual_Graph_is_a_Tree}
    \uses{lem:Existence_of_Leaves_in_a_Tree}
    \uses{lem:Dual_Graph_Leaf_is_Ear}
    The proof is by induction on $n$, the number of vertices. The base case $n=3$ is trivial. For $n>3$, consider the dual graph of the triangulation, which is a tree by lemma \ref{lem:Dual_Graph_is_a_Tree}. By lemma \ref{lem:Existence_of_Leaves_in_a_Tree}, this tree has a leaf. By lemma \ref{lem:Dual_Graph_Leaf_is_Ear}, this leaf corresponds to an ear triangle of the polygon. Let this ear be $\triangle v_1v_2v_3$, where $\overline{v_1v_3}$ is the diagonal.
    
    Removing this ear (and vertex $v_2$) leaves a smaller polygon with $n-1$ vertices. By the induction hypothesis, this smaller polygon is 3-colorable. In any such coloring, vertices $v_1$ and $v_3$ are adjacent and thus have different colors. We can then color the vertex $v_2$ with the third color not used by $v_1$ or $v_3$. This yields a valid 3-coloring for the entire polygon.
\end{proof}

\begin{lemma}[Vertices of a Triangle in a 3-Colored Triangulation]
    \label{lem:Vertices_of_a_Triangle_in_a_3_Colored_Triangulation}
    \uses{def:3-Coloring_of_a_Graph}
    \uses{def:Polygon_Triangulation}
    In any 3-coloring of the graph of a triangulated polygon, the three vertices of any triangle in the triangulation must have distinct colors.
\end{lemma}
\begin{proof}
    In a triangulation, the three vertices of any triangle are pairwise connected by edges (either edges of the original polygon or diagonals of the triangulation). By the definition of a 3-coloring, any two adjacent vertices must have different colors. Therefore, the three vertices of a triangle, being all mutually adjacent, must be assigned three distinct colors.
\end{proof}

\begin{lemma}[Existence of a Small Color Class]
    \label{lem:Existence_of_a_Small_Color_Class}
    For any 3-coloring of a set of $n$ vertices, there is at least one color class with at most $\lfloor n/3 \rfloor$ vertices.
\end{lemma}

\begin{proof}
    Let the sizes of the three color classes be $n_1, n_2,$ and $n_3$. We have $n_1 + n_2 + n_3 = n$. Assume for contradiction that $n_i > \lfloor n/3 \rfloor$ for all $i \in \{1, 2, 3\}$. Since $n_i$ must be an integer, this means $n_i \ge \lfloor n/3 \rfloor + 1$ for all $i$. Summing these inequalities gives $n_1 + n_2 + n_3 \ge 3(\lfloor n/3 \rfloor + 1)$. This sum is greater than $n$, which contradicts $n_1 + n_2 + n_3 = n$. Thus, at least one class must have size at most $\lfloor n/3 \rfloor$.
\end{proof}

\begin{lemma}[Guards on Vertices of a Triangle]
    \label{lem:Guards_on_Vertices_of_a_Triangle}
    A guard placed at any vertex of a triangle can see the entire triangle.
\end{lemma}

\begin{proof}
    A triangle is a convex set. By definition of convexity, for any two points $p, q$ in a triangle, the line segment $\overline{pq}$ is entirely contained within the triangle. Therefore, a guard placed at a vertex can see all other points in the triangle.
\end{proof}

\begin{theorem}[Chvátal's Art Gallery Theorem]
    \label{thm:Chvatals_Art_Gallery_Theorem}
    \uses{def:Guard_Set}
    For any simple polygon with $n$ vertices, $\lfloor n/3 \rfloor$ guards are sufficient to guard the entire polygon.
\end{theorem}

\begin{proof}
    \uses{lem:Existence_of_Triangulation}
    \uses{lem:3-Colorability_of_a_Triangulated_Polygon}
    \uses{lem:Existence_of_a_Small_Color_Class}
    \uses{lem:Vertices_of_a_Triangle_in_a_3_Colored_Triangulation}
    \uses{lem:Guards_on_Vertices_of_a_Triangle}
    First, we triangulate the simple polygon according to lemma \ref{lem:Existence_of_Triangulation}. Next, we apply a 3-coloring to the vertices of the triangulated polygon, which is guaranteed to exist by lemma \ref{lem:3-Colorability_of_a_Triangulated_Polygon}. This coloring partitions the $n$ vertices into three color classes.
    
    By lemma \ref{lem:Existence_of_a_Small_Color_Class}, there exists a color class with size at most $\lfloor n/3 \rfloor$. We place guards at all vertices belonging to this smallest color class.
    
    Now consider any triangle in the triangulation. By lemma \ref{lem:Vertices_of_a_Triangle_in_a_3_Colored_Triangulation}, its three vertices must have distinct colors. Therefore, every triangle contains exactly one vertex from our chosen color class, and thus has exactly one guard.
    
    By lemma \ref{lem:Guards_on_Vertices_of_a_Triangle}, this single guard is sufficient to see the entire triangle. Since the union of all triangles in the triangulation is the polygon itself, the entire polygon is guarded. Thus, $\lfloor n/3 \rfloor$ guards are sufficient.
\end{proof}

\chapter{Fermat's Theorem on Stationary Points}

\begin{definition}[Local Extremum]
    \label{def:Local_Extremum}
    A function $f$ has a local maximum at a point $c$ if there exists an open interval $I$ containing $c$ such that $f(x) \leq f(c)$ for all $x \in I$. Similarly, $f$ has a local minimum at $c$ if $f(x) \geq f(c)$ for all $x \in I$. A local extremum is either a local maximum or a local minimum.
\end{definition}

\begin{definition}[Stationary Point]
    \label{def:Stationary_Point}
    A point $c$ in the domain of a differentiable function $f$ is called a stationary point if $f'(c) = 0$.
\end{definition}

\begin{lemma}[Preservation of Inequality under Limits]
    \label{lem:Preservation_of_Inequality_under_Limits}
    Let $g$ be a real-valued function defined on an open interval containing a point $c$.
    If $g(x) \leq K$ for all $x$ in some open interval around $c$ (excluding possibly $c$ itself), and $\lim_{x \to c} g(x)$ exists, then $\lim_{x \to c} g(x) \leq K$.
    Similarly, if $g(x) \geq K$, then $\lim_{x \to c} g(x) \geq K$. This holds for one-sided limits as well.
\end{lemma}

\begin{proof}
    We prove the first statement by contradiction. Assume $\lim_{x \to c} g(x) = L > K$. Let $\epsilon = L - K > 0$. By the definition of a limit, there exists a $\delta > 0$ such that for all $x$ with $0 < |x-c| < \delta$, we have $|g(x) - L| < \epsilon$. This implies $L - \epsilon < g(x) < L + \epsilon$. Substituting $\epsilon = L - K$, we get $K < g(x)$. This contradicts the hypothesis that $g(x) \leq K$ for all $x$ in a neighborhood of $c$. Thus, our assumption must be false, and $\lim_{x \to c} g(x) \leq K$. The proof for $g(x) \geq K$ is analogous.
\end{proof}

\begin{lemma}[Existence of a Two-Sided Limit]
    \label{lem:Two_Sided_Limit_Existence}
    Let $g$ be a real-valued function defined on an open interval containing a point $c$, except possibly at $c$ itself. The limit $\lim_{x \to c} g(x)$ exists and equals $L$ if and only if both one-sided limits, $\lim_{x \to c^-} g(x)$ and $\lim_{x \to c^+} g(x)$, exist and are equal to $L$.
\end{lemma}

\begin{proof}
    ($\Rightarrow$) Assume $\lim_{x \to c} g(x) = L$. By the $\epsilon$-$\delta$ definition of a limit, for every $\epsilon > 0$, there exists a $\delta > 0$ such that if $0 < |x-c| < \delta$, then $|g(x)-L| < \epsilon$. This single $\delta$ works for both $x > c$ and $x < c$. Specifically, if $c < x < c+\delta$, then $|g(x)-L| < \epsilon$, which implies $\lim_{x \to c^+} g(x) = L$. Similarly, if $c-\delta < x < c$, then $|g(x)-L| < \epsilon$, which implies $\lim_{x \to c^-} g(x) = L$.
    ($\Leftarrow$) Assume $\lim_{x \to c^-} g(x) = L$ and $\lim_{x \to c^+} g(x) = L$. For any $\epsilon > 0$, there exists $\delta_1 > 0$ such that if $c-\delta_1 < x < c$, then $|g(x)-L| < \epsilon$. Also, there exists $\delta_2 > 0$ such that if $c < x < c+\delta_2$, then $|g(x)-L| < \epsilon$. Let $\delta = \min(\delta_1, \delta_2)$. Then for any $x$ such that $0 < |x-c| < \delta$, the condition $|g(x)-L| < \epsilon$ is satisfied. This is precisely the definition of $\lim_{x \to c} g(x) = L$.
\end{proof}

\begin{lemma}[Condition for Differentiability]
    \label{lem:Condition_for_Differentiability}
    A function $f$ is differentiable at a point $c$ if and only if its left-hand derivative $f'_-(c) = \lim_{h \to 0^-} \frac{f(c+h)-f(c)}{h}$ and its right-hand derivative $f'_+(c) = \lim_{h \to 0^+} \frac{f(c+h)-f(c)}{h}$ both exist and are equal. In this case, $f'(c) = f'_-(c) = f'_+(c)$.
\end{lemma}

\begin{proof}
    \uses{lem:Two_Sided_Limit_Existence}
    The derivative of $f$ at $c$ is defined as the limit $f'(c) = \lim_{h \to 0} \frac{f(c+h)-f(c)}{h}$. Let $g(h) = \frac{f(c+h)-f(c)}{h}$. According to lemma \ref{lem:Two_Sided_Limit_Existence}, the two-sided limit $\lim_{h \to 0} g(h)$ exists if and only if the one-sided limits $\lim_{h \to 0^-} g(h)$ (the left-hand derivative) and $\lim_{h \to 0^+} g(h)$ (the right-hand derivative) exist and are equal.
\end{proof}

\begin{lemma}[Trichotomy Property Consequence]
    \label{lem:Trichotomy_Property_Consequence}
    For any real number $x$, if $x \leq 0$ and $x \geq 0$, then $x=0$.
\end{lemma}

\begin{proof}
    This follows from the trichotomy property of real numbers, which states that for any real number $x$, exactly one of the following is true: $x < 0$, $x > 0$, or $x = 0$. The conditions $x \leq 0$ and $x \geq 0$ rule out $x < 0$ and $x > 0$, respectively. Therefore, the only remaining possibility is $x=0$.
\end{proof}

\begin{lemma}[One-Sided Derivatives at a Local Maximum]
    \label{lem:One_Sided_Derivatives_at_a_Local_Maximum}
    \uses{def:Local_Extremum}
    If a function $f$ has a local maximum at a point $c$, and the one-sided derivatives at $c$ exist, then the right-hand derivative $f'_+(c) \leq 0$ and the left-hand derivative $f'_-(c) \geq 0$.
\end{lemma}

\begin{proof}
    \uses{def:Local_Extremum}
    \uses{lem:Preservation_of_Inequality_under_Limits}
    Since $f$ has a local maximum at $c$, there exists an open interval around $c$ such that for any $x$ in this interval, $f(x) \leq f(c)$.
    For $h > 0$ small enough such that $c+h$ is in this interval, we have $f(c+h) \leq f(c)$, which implies $f(c+h) - f(c) \leq 0$. Since $h > 0$, the difference quotient is $\frac{f(c+h)-f(c)}{h} \leq 0$. Applying lemma \ref{lem:Preservation_of_Inequality_under_Limits}, we have $f'_+(c) = \lim_{h \to 0^+} \frac{f(c+h)-f(c)}{h} \leq 0$.
    For $h < 0$ small enough, we have $f(c+h) \leq f(c)$, so $f(c+h) - f(c) \leq 0$. Since $h < 0$, the difference quotient is $\frac{f(c+h)-f(c)}{h} \geq 0$. Applying lemma \ref{lem:Preservation_of_Inequality_under_Limits}, we have $f'_-(c) = \lim_{h \to 0^-} \frac{f(c+h)-f(c)}{h} \geq 0$.
\end{proof}

\begin{lemma}[One-Sided Derivatives at a Local Minimum]
    \label{lem:One_Sided_Derivatives_at_a_Local_Minimum}
    \uses{def:Local_Extremum}
    If a function $f$ has a local minimum at a point $c$, and the one-sided derivatives at $c$ exist, then the right-hand derivative $f'_+(c) \geq 0$ and the left-hand derivative $f'_-(c) \leq 0$.
\end{lemma}

\begin{proof}
    \uses{def:Local_Extremum}
    \uses{lem:Preservation_of_Inequality_under_Limits}
    Since $f$ has a local minimum at $c$, for $h > 0$ small enough such that $c+h$ is in the neighborhood of $c$, we have $f(c+h) \geq f(c)$, which implies $f(c+h) - f(c) \geq 0$. As $h > 0$, the difference quotient $\frac{f(c+h)-f(c)}{h} \geq 0$. By lemma \ref{lem:Preservation_of_Inequality_under_Limits}, $f'_+(c) = \lim_{h \to 0^+} \frac{f(c+h)-f(c)}{h} \geq 0$.
    For $h < 0$ small enough, we have $f(c+h) \geq f(c)$, so $f(c+h) - f(c) \geq 0$. As $h < 0$, the difference quotient $\frac{f(c+h)-f(c)}{h} \leq 0$. By lemma \ref{lem:Preservation_of_Inequality_under_Limits}, $f'_-(c) = \lim_{h \to 0^-} \frac{f(c+h)-f(c)}{h} \leq 0$.
\end{proof}

\begin{theorem}[Fermat's Theorem on Stationary Points]
    \label{thm:Fermats_Theorem_on_Stationary_Points}
    \uses{def:Local_Extremum}
    \uses{def:Stationary_Point}
    If a function $f$ has a local extremum at a point $c$ in an open interval, and if $f$ is differentiable at $c$, then $c$ is a stationary point of $f$, i.e., $f'(c) = 0$.
\end{theorem}

\begin{proof}
    \uses{lem:Condition_for_Differentiability}
    \uses{lem:Trichotomy_Property_Consequence}
    \uses{lem:One_Sided_Derivatives_at_a_Local_Maximum}
    \uses{lem:One_Sided_Derivatives_at_a_Local_Minimum}
    Case 1: $f$ has a local maximum at $c$.
    Since $f$ is differentiable, its one-sided derivatives exist. By lemma \ref{lem:One_Sided_Derivatives_at_a_Local_Maximum}, we know $f'_+(c) \leq 0$ and $f'_-(c) \geq 0$.
    By lemma \ref{lem:Condition_for_Differentiability}, differentiability at $c$ implies $f'(c) = f'_+(c) = f'_-(c)$.
    Therefore, we have $f'(c) \leq 0$ and $f'(c) \geq 0$. By lemma \ref{lem:Trichotomy_Property_Consequence}, this implies $f'(c) = 0$.

    Case 2: $f$ has a local minimum at $c$.
    Since $f$ is differentiable, its one-sided derivatives exist. By lemma \ref{lem:One_Sided_Derivatives_at_a_Local_Minimum}, we know $f'_+(c) \geq 0$ and $f'_-(c) \leq 0$.
    By lemma \ref{lem:Condition_for_Differentiability}, $f'(c) = f'_+(c) = f'_-(c)$.
    Therefore, we have $f'(c) \geq 0$ and $f'(c) \leq 0$. By lemma \ref{lem:Trichotomy_Property_Consequence}, this implies $f'(c) = 0$.

    In both cases, $f'(c)=0$, which means $c$ is a stationary point.
\end{proof}

\input{theorems/gauss-lucas}

\chapter{Hinge Theorem}

\begin{definition}[Law of Cosines]
    \label{def:Law_of_Cosines}
    For a triangle with sides of length $a$, $b$, and $c$, and with $\gamma$ as the angle opposite the side of length $c$, the Law of Cosines states that:
    $$c^2 = a^2 + b^2 - 2ab \cos \gamma$$
\end{definition}

\begin{lemma}[Triangle Angle Sum Theorem]
    \label{lem:Triangle_Angle_Sum}
    In Euclidean geometry, the sum of the measures of the interior angles of any non-degenerate triangle is equal to $\pi$ radians.
\end{lemma}

\begin{proof}
    This is a fundamental theorem in Euclidean geometry, derived from the parallel postulate. A formal proof involves constructing a line parallel to one side of the triangle through the opposite vertex and using properties of alternate interior angles.
\end{proof}

\begin{lemma}[Mean Value Theorem]
    \label{lem:Mean_Value_Theorem}
    Let $f$ be a function that is continuous on a closed interval $[a, b]$ and differentiable on the open interval $(a, b)$. Then there exists at least one point $c \in (a, b)$ such that $f'(c) = \frac{f(b) - f(a)}{b - a}$.
\end{lemma}

\begin{proof}
    This is a fundamental theorem of differential calculus. Its proof typically relies on Rolle's Theorem, which is a special case of the Mean Value Theorem where $f(a) = f(b)$.
\end{proof}

\begin{lemma}[Derivative of Cosine]
    \label{lem:Derivative_of_Cosine}
    The function $f(\gamma) = \cos \gamma$ is differentiable for all $\gamma \in \mathbb{R}$, and its derivative is given by $f'(\gamma) = -\sin \gamma$.
\end{lemma}

\begin{proof}
    This is a standard result from differential calculus, typically proven from the limit definition of the derivative, using the angle addition formula for cosine and the fundamental trigonometric limits $\lim_{h\to 0} \frac{\sin h}{h} = 1$ and $\lim_{h\to 0} \frac{\cos h - 1}{h} = 0$.
\end{proof}

\begin{lemma}[Inequality Property of Multiplication by a Negative Number]
    \label{lem:Inequality_Multiplication_Property}
    Let $x, y, z$ be real numbers. If $x < y$ and $z < 0$, then $zx > zy$.
\end{lemma}

\begin{proof}
    Given $x < y$, it follows that $y - x > 0$. Given $z < 0$, it follows that $-z > 0$. The product of two positive real numbers is positive, so $(y - x)(-z) > 0$. Distributing $-z$ gives $-yz + xz > 0$, which is equivalent to $xz > yz$.
\end{proof}

\begin{lemma}[Properties of Triangle Angles]
    \label{lem:Triangle_Angle_Properties}
    \uses{lem:Triangle_Angle_Sum}
    The measure of any interior angle of a non-degenerate triangle is a real number in the interval $(0, \pi)$.
\end{lemma}

\begin{proof}
    \uses{lem:Triangle_Angle_Sum}
    Let the interior angles of a triangle be $\alpha, \beta, \gamma$. By lemma \ref{lem:Triangle_Angle_Sum}, we have $\alpha + \beta + \gamma = \pi$. In a non-degenerate triangle, the measure of each angle must be positive, so $\alpha > 0$, $\beta > 0$, and $\gamma > 0$. From the sum, we can write $\gamma = \pi - (\alpha + \beta)$. Since $\alpha > 0$ and $\beta > 0$, their sum $(\alpha + \beta)$ is positive. Therefore, $\gamma < \pi$. Combining these results, we have $0 < \gamma < \pi$. The same argument applies to $\alpha$ and $\beta$.
\end{proof}

\begin{lemma}[Strict Monotonicity and Derivatives]
    \label{lem:Monotonicity_and_Derivatives}
    \uses{lem:Mean_Value_Theorem}
    Let $f$ be a real-valued function that is continuous on an interval $[a, b]$ and differentiable on $(a, b)$. If $f'(x) < 0$ for all $x \in (a, b)$, then $f$ is strictly decreasing on $[a, b]$.
\end{lemma}

\begin{proof}
    \uses{lem:Mean_Value_Theorem}
    Let $x_1, x_2 \in [a, b]$ with $x_1 < x_2$. By lemma \ref{lem:Mean_Value_Theorem} applied to the interval $[x_1, x_2]$, there exists a point $c \in (x_1, x_2)$ such that $f(x_2) - f(x_1) = f'(c)(x_2 - x_1)$. By hypothesis, $f'(c) < 0$. Since $x_1 < x_2$, we have $(x_2 - x_1) > 0$. The product of a negative number and a positive number is negative, so $f'(c)(x_2 - x_1) < 0$. This implies $f(x_2) - f(x_1) < 0$, or $f(x_2) < f(x_1)$. Since this holds for any $x_1 < x_2$ in the interval, the function $f$ is strictly decreasing on $[a, b]$.
\end{proof}

\begin{lemma}[Positivity of Sine on (0, pi)]
    \label{lem:Positivity_of_Sine}
    For any angle $\gamma \in (0, \pi)$, $\sin \gamma > 0$.
\end{lemma}

\begin{proof}
    From the unit circle definition of trigonometric functions, $\sin \gamma$ corresponds to the y-coordinate of a point on the circle. For angles $\gamma$ in the interval $(0, \pi)$, the point lies in the first or second quadrant, where the y-coordinate is positive.
\end{proof}

\begin{lemma}[Order Preservation for Positive Numbers under Squaring]
    \label{lem:Order_Preservation_Squaring}
    Let $c$ and $c'$ be positive real numbers. Then $c^2 > c'^2$ if and only if $c > c'$.
\end{lemma}

\begin{proof}
    The inequality $c^2 > c'^2$ is equivalent to $c^2 - c'^2 > 0$. Factoring the left side gives $(c - c')(c + c') > 0$. Since $c$ and $c'$ are positive real numbers, their sum $(c + c')$ is positive. For the product of two factors to be positive, if one factor is positive, the other must also be positive. Thus, we must have $(c - c') > 0$, which implies $c > c'$. Since all steps are reversible, the equivalence holds.
\end{proof}

\begin{lemma}[Monotonicity of the Cosine Function]
    \label{lem:Monotonicity_of_Cosine}
    \uses{lem:Derivative_of_Cosine}
    \uses{lem:Monotonicity_and_Derivatives}
    \uses{lem:Positivity_of_Sine}
    The cosine function is a strictly decreasing function on the interval $[0, \pi]$.
\end{lemma}

\begin{proof}
    \uses{lem:Derivative_of_Cosine}
    \uses{lem:Monotonicity_and_Derivatives}
    \uses{lem:Positivity_of_Sine}
    Let $f(\gamma) = \cos \gamma$. This function is continuous on $[0, \pi]$ and differentiable on $(0, \pi)$. By lemma \ref{lem:Derivative_of_Cosine}, its derivative is $f'(\gamma) = -\sin \gamma$. By lemma \ref{lem:Positivity_of_Sine}, for any $\gamma \in (0, \pi)$, we have $\sin \gamma > 0$. Consequently, $f'(\gamma) = -\sin \gamma < 0$ for all $\gamma \in (0, \pi)$. By lemma \ref{lem:Monotonicity_and_Derivatives}, the function $\cos \gamma$ is strictly decreasing on the interval $[0, \pi]$.
\end{proof}

\begin{theorem}[Hinge Theorem]
    \label{thm:Hinge_Theorem}
    \uses{def:Law_of_Cosines}
    \uses{lem:Triangle_Angle_Properties}
    \uses{lem:Inequality_Multiplication_Property}
    \uses{lem:Order_Preservation_Squaring}
    \uses{lem:Monotonicity_of_Cosine}
    Let $\triangle ABC$ and $\triangle A'B'C'$ be two triangles. If $AC = A'C' = b$, $BC = B'C' = a$, and $\angle C > \angle C'$, then the length of the third side $AB$ is greater than the length of the third side $A'B'$.
\end{theorem}

\begin{proof}
    \uses{def:Law_of_Cosines}
    \uses{lem:Triangle_Angle_Properties}
    \uses{lem:Monotonicity_of_Cosine}
    \uses{lem:Inequality_Multiplication_Property}
    \uses{lem:Order_Preservation_Squaring}
    Let $c = AB$, $c' = A'B'$, $\gamma = \angle C$, and $\gamma' = \angle C'$. By applying the Law of Cosines from definition \ref{def:Law_of_Cosines} to both triangles, we have:
    $$c^2 = a^2 + b^2 - 2ab \cos \gamma$$
    $$c'^2 = a^2 + b^2 - 2ab \cos \gamma'$$
    By lemma \ref{lem:Triangle_Angle_Properties}, both $\gamma$ and $\gamma'$ are in the interval $(0, \pi)$. We are given that $\gamma > \gamma'$. By lemma \ref{lem:Monotonicity_of_Cosine}, the cosine function is strictly decreasing on $[0, \pi]$, so $\cos \gamma < \cos \gamma'$.
    Since $a$ and $b$ are side lengths, they are positive, so $-2ab < 0$. Using lemma \ref{lem:Inequality_Multiplication_Property} with $x = \cos \gamma$, $y = \cos \gamma'$, and $z = -2ab$, we can multiply the inequality by $-2ab$ and reverse the sign:
    $$-2ab \cos \gamma > -2ab \cos \gamma'$$
    Adding $a^2 + b^2$ to both sides yields:
    $$a^2 + b^2 - 2ab \cos \gamma > a^2 + b^2 - 2ab \cos \gamma'$$
    This is equivalent to $c^2 > c'^2$. Since $c$ and $c'$ are positive side lengths, lemma \ref{lem:Order_Preservation_Squaring} implies that $c > c'$. Thus, $AB > A'B'$.
\end{proof}

\chapter{Hurwitz's Theorem}

\begin{definition}[Holomorphic Function]
    \label{def:Holomorphic_Function}
    A function $f: U \to \mathbb{C}$, where $U$ is an open subset of the complex plane $\mathbb{C}$, is said to be \textit{holomorphic} on $U$ if it is complex differentiable at every point $z_0 \in U$.
\end{definition}

\begin{definition}[Uniform Convergence on Compact Subsets]
    \label{def:Uniform_Convergence_on_Compact_Subsets}
    \uses{def:Holomorphic_Function}
    A sequence of functions $\{f_n\}$ is said to converge uniformly on compact subsets of a domain $G$ to a function $f$ if for every compact set $K \subset G$, the sequence of functions $\{f_n|_K\}$ converges uniformly to $f|_K$.
\end{definition}

\begin{lemma}[Green's Theorem]
    \label{lem:Greens_Theorem}
    Let $C$ be a positively oriented, piecewise smooth, simple closed curve in a plane, and let $D$ be the region bounded by $C$. If $L(x,y)$ and $M(x,y)$ have continuous first partial derivatives on an open region containing $D$, then
    $$ \oint_C (L \, dx + M \, dy) = \iint_D \left(\frac{\partial M}{\partial x} - \frac{\partial L}{\partial y}\right) dA. $$
\end{lemma}

\begin{proof}
    This is a fundamental theorem of vector calculus. The proof for a simple rectangular region is straightforward. For a general region $D$, it can be approximated by a grid of small rectangles, and the line integrals over the interior boundaries cancel out, leaving the integral over the boundary $C$.
\end{proof}

\begin{lemma}[Cauchy's Integral Theorem]
    \label{lem:Cauchy_Integral_Theorem}
    \uses{def:Holomorphic_Function}
    \uses{lem:Greens_Theorem}
    Let $U$ be a simply connected open subset of $\mathbb{C}$, let $f: U \to \mathbb{C}$ be a holomorphic function, and let $\gamma$ be a simple closed contour in $U$. Then $\oint_\gamma f(z) dz = 0$.
\end{lemma}

\begin{proof}
    Let $f(z) = u(x,y) + i v(x,y)$ and $dz = dx + i dy$. The integral is $\oint_\gamma (u+iv)(dx+idy) = \oint_\gamma (u \, dx - v \, dy) + i \oint_\gamma (v \, dx + u \, dy)$.
    Applying Green's Theorem, lemma \ref{lem:Greens_Theorem}, to the real and imaginary parts, we get:
    $$ \oint_\gamma (u \, dx - v \, dy) = \iint_D \left(\frac{\partial(-v)}{\partial x} - \frac{\partial u}{\partial y}\right) dA = \iint_D \left(-\frac{\partial v}{\partial x} - \frac{\partial u}{\partial y}\right) dA $$
    $$ \oint_\gamma (v \, dx + u \, dy) = \iint_D \left(\frac{\partial u}{\partial x} - \frac{\partial v}{\partial y}\right) dA $$
    Since $f$ is holomorphic, its real and imaginary parts $u$ and $v$ satisfy the Cauchy-Riemann equations: $\frac{\partial u}{\partial x} = \frac{\partial v}{\partial y}$ and $\frac{\partial u}{\partial y} = -\frac{\partial v}{\partial x}$. Substituting these equations into the double integrals shows that both integrands are zero. Therefore, the integral $\oint_\gamma f(z) dz$ is zero.
\end{proof}

\begin{lemma}[Deformation Invariance]
    \label{lem:Deformation_Invariance}
    \uses{def:Holomorphic_Function}
    \uses{lem:Cauchy_Integral_Theorem}
    Let $f$ be a holomorphic function in a domain $G$. If $\gamma_0$ and $\gamma_1$ are two simple closed contours in $G$ that are homotopic (can be continuously deformed into one another) within $G$, then $\oint_{\gamma_0} f(z) dz = \oint_{\gamma_1} f(z) dz$.
\end{lemma}

\begin{proof}
    Consider the closed contour $\Gamma = \gamma_0 - \gamma_1$ (where $-\gamma_1$ denotes the curve $\gamma_1$ with opposite orientation). Since $\gamma_0$ and $\gamma_1$ are homotopic in $G$, the region enclosed by $\Gamma$ is contained in $G$. We can construct a contour that connects $\gamma_0$ and $\gamma_1$, traverses $\gamma_0$, goes back along the connection, traverses $\gamma_1$ in reverse, and returns along the connection. The integrals along the connecting path cancel. The region enclosed by this new composite contour is simply connected and lies within $G$. By Cauchy's Integral Theorem, lemma \ref{lem:Cauchy_Integral_Theorem}, the integral of $f$ over this composite contour is zero. This implies $\oint_{\gamma_0} f(z) dz - \oint_{\gamma_1} f(z) dz = 0$.
\end{proof}

\begin{lemma}[Cauchy's Integral Formula for a Simple Pole]
    \label{lem:Cauchy_Integral_Formula_for_a_Simple_Pole}
    \uses{lem:Cauchy_Integral_Theorem}
    \uses{lem:Deformation_Invariance}
    Let $\gamma$ be a simple closed contour traversed counter-clockwise, and let $a$ be a point not on $\gamma$. Then
    $$ \oint_\gamma \frac{dz}{z-a} = \begin{cases} 2\pi i & \text{if } a \text{ is inside } \gamma \\ 0 & \text{if } a \text{ is outside } \gamma \end{cases}. $$
\end{lemma}

\begin{proof}
    If $a$ is outside $\gamma$, the function $f(z) = 1/(z-a)$ is holomorphic inside and on $\gamma$. By Cauchy's Integral Theorem, lemma \ref{lem:Cauchy_Integral_Theorem}, the integral is 0.
    If $a$ is inside $\gamma$, let $C_r$ be a small circle of radius $r$ centered at $a$, entirely inside $\gamma$. The contours $\gamma$ and $C_r$ are homotopic in the domain $\mathbb{C} \setminus \{a\}$. By the Deformation Invariance, lemma \ref{lem:Deformation_Invariance}, $\oint_\gamma \frac{dz}{z-a} = \oint_{C_r} \frac{dz}{z-a}$. Parametrizing $C_r$ by $z(t) = a + re^{it}$ for $t \in [0, 2\pi]$, we have $dz = ire^{it}dt$. The integral becomes $\int_0^{2\pi} \frac{ire^{it}}{re^{it}} dt = \int_0^{2\pi} i dt = 2\pi i$.
\end{proof}

\begin{lemma}[Residue of Logarithmic Derivative]
    \label{lem:Residue_of_Logarithmic_Derivative}
    \uses{def:Holomorphic_Function}
    Let $f$ be a function analytic at $z_0$. If $f$ has a zero of order $m$ at $z_0$, the logarithmic derivative $f'(z)/f(z)$ has a simple pole with residue $m$ at $z_0$. If $f$ has a pole of order $n$ at $z_0$, the residue is $-n$.
\end{lemma}

\begin{proof}
    If $f$ has a zero of order $m$ at $z_0$, we can write $f(z)=(z-z_0)^m g(z)$ where $g(z)$ is analytic and non-zero at $z_0$. Taking the derivative, $f'(z) = m(z-z_0)^{m-1}g(z) + (z-z_0)^m g'(z)$. The logarithmic derivative is $\frac{f'(z)}{f(z)} = \frac{m(z-z_0)^{m-1}g(z) + (z-z_0)^m g'(z)}{(z-z_0)^m g(z)} = \frac{m}{z-z_0} + \frac{g'(z)}{g(z)}$. Since $g(z_0)\neq 0$, the term $g'(z)/g(z)$ is analytic at $z_0$. Thus, $f'(z)/f(z)$ has a simple pole at $z_0$ with residue $m$.
    If $f$ has a pole of order $n$, $f(z)=(z-z_0)^{-n} h(z)$ with $h(z_0)\neq 0$. A similar calculation yields $\frac{f'(z)}{f(z)} = \frac{-n}{z-z_0} + \frac{h'(z)}{h(z)}$, showing the residue is $-n$.
\end{proof}

\begin{lemma}[Argument Principle]
    \label{lem:Argument_Principle}
    \uses{def:Holomorphic_Function}
    \uses{lem:Cauchy_Integral_Theorem}
    \uses{lem:Cauchy_Integral_Formula_for_a_Simple_Pole}
    \uses{lem:Residue_of_Logarithmic_Derivative}
    Let $f$ be a meromorphic function on a domain $G$, and let $\gamma$ be a simple closed contour in $G$ that does not pass through any zeros or poles of $f$. Let $Z$ be the number of zeros and $P$ be the number of poles of $f$ inside $\gamma$, counted with multiplicity. Then
    $$ \frac{1}{2\pi i} \oint_\gamma \frac{f'(z)}{f(z)} dz = Z - P. $$
\end{lemma}

\begin{proof}
    Let the zeros be $z_j$ with multiplicity $m_j$ and poles be $p_k$ with multiplicity $n_k$. By lemma \ref{lem:Residue_of_Logarithmic_Derivative}, the function $h(z) = f'(z)/f(z)$ has simple poles at each $z_j$ and $p_k$ with residues $m_j$ and $-n_k$ respectively. We can write $h(z) = \sum_j \frac{m_j}{z-z_j} - \sum_k \frac{n_k}{z-p_k} + g(z)$, where $g(z)$ is holomorphic inside $\gamma$. Integrating over $\gamma$ gives $\oint_\gamma h(z) dz = \sum_j m_j \oint_\gamma \frac{dz}{z-z_j} - \sum_k n_k \oint_\gamma \frac{dz}{z-p_k} + \oint_\gamma g(z) dz$. By lemma \ref{lem:Cauchy_Integral_Formula_for_a_Simple_Pole}, the integrals involving zeros and poles evaluate to $2\pi i$. By lemma \ref{lem:Cauchy_Integral_Theorem}, the integral of $g(z)$ is zero. Thus, $\oint_\gamma \frac{f'(z)}{f(z)} dz = 2\pi i (\sum_j m_j - \sum_k n_k) = 2\pi i (Z - P)$.
\end{proof}

\begin{lemma}[Taylor Series Expansion of Holomorphic Functions]
    \label{lem:Taylor_Series_Expansion}
    \uses{def:Holomorphic_Function}
    If a function $f$ is holomorphic in a disk $D(z_0, R)$, then $f$ can be represented by a power series $f(z) = \sum_{n=0}^\infty a_n (z-z_0)^n$ for all $z \in D(z_0, R)$. If $f$ is not the zero function, there exists a smallest integer $m \ge 0$ such that $a_m \neq 0$.
\end{lemma}

\begin{proof}
    The coefficients of the Taylor series are given by Cauchy's integral formula for derivatives: $a_n = \frac{f^{(n)}(z_0)}{n!} = \frac{1}{2\pi i} \oint_C \frac{f(w)}{(w-z_0)^{n+1}} dw$. The convergence of this series to $f(z)$ within the disk is a standard result in complex analysis. If $f$ is not identically zero, not all coefficients can be zero. Thus, there must be a smallest index $m$ for which $a_m \neq 0$.
\end{proof}

\begin{lemma}[Isolation of Zeros]
    \label{lem:Isolation_of_Zeros}
    \uses{def:Holomorphic_Function}
    \uses{lem:Taylor_Series_Expansion}
    If $f$ is a non-zero holomorphic function on a connected open set $G \subset \mathbb{C}$, then its zeros are isolated.
\end{lemma}

\begin{proof}
    Let $z_0$ be a zero of $f$. Since $f$ is not identically zero, by lemma \ref{lem:Taylor_Series_Expansion}, there is a Taylor expansion $f(z) = \sum_{k=m}^\infty a_k (z-z_0)^k$ with $m \ge 1$ and $a_m \neq 0$. We can factor this as $f(z) = (z-z_0)^m \left(a_m + a_{m+1}(z-z_0) + \dots\right) = (z-z_0)^m g(z)$. The function $g(z)$ is holomorphic in a neighborhood of $z_0$ and $g(z_0) = a_m \neq 0$. Since $g$ is continuous, there is a disk $D(z_0, r)$ where $g(z) \neq 0$. In this disk, for $z \neq z_0$, $f(z) = (z-z_0)^m g(z) \neq 0$. Thus, $z_0$ is an isolated zero.
\end{proof}

\begin{lemma}[Extreme Value Theorem]
    \label{lem:Extreme_Value_Theorem}
    If $K$ is a compact topological space and $h: K \to \mathbb{R}$ is a continuous function, then $h$ attains its maximum and minimum values on $K$.
\end{lemma}

\begin{proof}
    Let $S = h(K)$ be the image of $K$. Since $h$ is continuous and $K$ is compact, $S$ is a compact subset of $\mathbb{R}$. As a compact subset of $\mathbb{R}$, $S$ is closed and bounded. Being bounded, it has a supremum and an infimum. Being closed, it contains its supremum and infimum. Therefore, there exist points $x_{\mathrm{max}}, x_{\mathrm{min}} \in K$ such that $h(x_{\mathrm{max}}) = \sup S$ and $h(x_{\mathrm{min}}) = \inf S$.
\end{proof}

\begin{lemma}[Minimum Modulus on a Compact Set]
    \label{lem:Minimum_Modulus_on_Compact_Set}
    \uses{lem:Extreme_Value_Theorem}
    Let $K \subset \mathbb{C}$ be a compact set and let $g: K \to \mathbb{C}$ be a continuous function. If $g(z) \neq 0$ for all $z \in K$, then there exists a constant $\delta > 0$ such that $|g(z)| \ge \delta$ for all $z \in K$.
\end{lemma}

\begin{proof}
    The function $|g(z)|$ is a continuous real-valued function on the compact set $K$. By the Extreme Value Theorem, lemma \ref{lem:Extreme_Value_Theorem}, $|g(z)|$ attains its minimum value on $K$. Let this minimum be $\delta = \min_{z \in K} |g(z)|$. Since $g(z) \neq 0$ for all $z \in K$, we have $|g(z)| > 0$ for all $z \in K$. Therefore, the minimum value $\delta$ must be strictly positive.
\end{proof}

\begin{lemma}[Continuity of Parametric Integral]
    \label{lem:Continuity_of_Parametric_Integral}
    Let $\gamma$ be a contour and let $F(t,z)$ be a function that is continuous in both variables $t \in [a,b]$ and $z \in \gamma$. Then the function $I(t) = \oint_\gamma F(t,z) dz$ is a continuous function of $t$ on $[a,b]$.
\end{lemma}

\begin{proof}
    Since the domain $[a,b] \times \gamma$ is compact, the continuous function $F(t,z)$ is uniformly continuous. Thus, for any $\epsilon > 0$, there exists a $\delta > 0$ such that if $|t_1 - t_2| < \delta$, then $|F(t_1, z) - F(t_2, z)| < \frac{\epsilon}{\mathrm{length}(\gamma)}$ for all $z \in \gamma$. Then $|I(t_1) - I(t_2)| = \left|\oint_\gamma (F(t_1,z) - F(t_2,z)) dz\right| \le \oint_\gamma |F(t_1,z) - F(t_2,z)| |dz| < \frac{\epsilon}{\mathrm{length}(\gamma)} \mathrm{length}(\gamma) = \epsilon$. This shows $I(t)$ is continuous.
\end{proof}

\begin{lemma}[Continuity of Integer-Valued Function]
    \label{lem:Continuity_of_Integer_Valued_Function}
    Let $f: X \to Y$ be a continuous function where $X$ is a connected topological space and $Y$ is a discrete topological space (e.g., $\mathbb{Z}$). Then $f$ must be a constant function.
\end{lemma}

\begin{proof}
    The image $f(X)$ of a connected space under a continuous map is connected. The only connected subsets of a discrete space like $\mathbb{Z}$ are single points. Therefore, $f(X)$ must be a single point, which means $f$ is constant.
\end{proof}

\begin{lemma}[Rouché's Theorem]
    \label{lem:Rouche's_Theorem}
    \uses{def:Holomorphic_Function}
    \uses{lem:Argument_Principle}
    \uses{lem:Continuity_of_Parametric_Integral}
    \uses{lem:Continuity_of_Integer_Valued_Function}
    Let $f$ and $g$ be holomorphic functions on a domain $G$. Let $\gamma$ be a simple closed contour in $G$, and suppose that for all $z \in \gamma$, we have $|f(z) - g(z)| < |f(z)|$. Then $f$ and $g$ have the same number of zeros inside $\gamma$.
\end{lemma}

\begin{proof}
    The condition $|f(z) - g(z)| < |f(z)|$ implies $f(z) \neq 0$ and $g(z) \neq 0$ on $\gamma$. Let $h(t, z) = f(z) + t(g(z)-f(z))$ for $t \in [0, 1]$. For any $t \in [0, 1]$ and $z \in \gamma$, $|h(t, z) - f(z)| = t|g(z)-f(z)| < t|f(z)| \le |f(z)|$, which implies $h(t,z) \neq 0$ on $\gamma$.
    Let $N(t)$ be the number of zeros of $h(t, \cdot)$ inside $\gamma$. By lemma \ref{lem:Argument_Principle}, $N(t) = \frac{1}{2\pi i} \oint_\gamma \frac{1}{h(t,z)}\frac{\partial h}{\partial z}(t,z) dz$. The integrand is continuous in $(t,z)$, so by lemma \ref{lem:Continuity_of_Parametric_Integral}, $N(t)$ is a continuous function of $t$. Since $N(t)$ must be an integer, lemma \ref{lem:Continuity_of_Integer_Valued_Function} implies $N(t)$ is constant on the connected interval $[0,1]$. Thus, $N(0) = N(1)$. $N(0)$ is the number of zeros of $h(0,z)=f(z)$, and $N(1)$ is the number of zeros of $h(1,z)=g(z)$. Hence, they are equal.
\end{proof}

\begin{theorem}[Hurwitz's Theorem]
    \label{thm:Hurwitzs_Theorem}
    \uses{def:Holomorphic_Function}
    \uses{def:Uniform_Convergence_on_Compact_Subsets}
    \uses{lem:Isolation_of_Zeros}
    \uses{lem:Minimum_Modulus_on_Compact_Set}
    \uses{lem:Rouche's_Theorem}
    Let $\{f_n\}$ be a sequence of holomorphic functions on a connected open set $G$ that converge uniformly on compact subsets of $G$ to a holomorphic function $f$ which is not constantly zero on $G$. If $f$ has a zero of order $m$ at $z_0$, then for every small enough $\rho > 0$ and for sufficiently large $k \in \mathbb{N}$, $f_k$ has precisely $m$ zeros in the disk defined by $|z - z_0| < \rho$, including multiplicity.
\end{theorem}

\begin{proof}
    Let $z_0$ be a zero of order $m$ of $f$. Since $f$ is not identically zero, lemma \ref{lem:Isolation_of_Zeros} ensures that we can choose $\rho > 0$ such that $f(z) \neq 0$ for $0 < |z-z_0| \le \rho$. Let $C$ be the circle $|z-z_0| = \rho$. $C$ is a compact set.
    As $f(z) \neq 0$ on $C$, by lemma \ref{lem:Minimum_Modulus_on_Compact_Set}, there exists a $\delta > 0$ such that $|f(z)| \ge \delta$ for all $z \in C$.
    The sequence $\{f_n\}$ converges to $f$ uniformly on $C$. Thus, there is an integer $N$ such that for all $n \ge N$ and $z \in C$, we have $|f_n(z) - f(z)| < \delta$.
    This gives the condition $|f_n(z) - f(z)| < \delta \le |f(z)|$ for all $n \ge N$ and $z \in C$.
    By Rouché's Theorem, lemma \ref{lem:Rouche's_Theorem}, the functions $f_n$ and $f$ have the same number of zeros inside the circle $C$ for all $n \ge N$.
    Since $f$ has a zero of order $m$ at $z_0$ and no other zeros within this disk, $f$ has exactly $m$ zeros inside $C$. Therefore, for all $n \ge N$, $f_n$ also has exactly $m$ zeros inside the disk $|z-z_0|<\rho$.
\end{proof}

\chapter{Kolmogorov-Arnold Representation Theorem}

\begin{definition}[Continuous Functions on the Unit Hypercube]
    \label{def:Continuous_Functions}
    Let $I = [0, 1]$. The space of all continuous functions from the $n$-dimensional unit hypercube $I^n$ to the real numbers $\mathbb{R}$ is denoted by $C(I^n)$. This is a normed vector space equipped with the supremum norm $\|f\| = \sup_{\mathbf{x} \in I^n} |f(\mathbf{x})|$.
\end{definition}

\begin{lemma}[Bolzano-Weierstrass Theorem for Sequence Compactness]
    \label{lem:Bolzano_Weierstrass}
    A set $K \subset \mathbb{R}^n$ is compact if and only if every sequence in $K$ has a subsequence that converges to a point in $K$.
\end{lemma}

\begin{proof}
    We prove the "if" direction for a closed and bounded set $K=I^n$. Let $\{\mathbf{x}_k\}$ be a sequence in $I^n$. Since $I^n$ is bounded, the sequence is bounded. We apply the bisection method. Let $C_0 = I^n$. Divide $C_0$ into $2^n$ closed sub-hypercubes. At least one, call it $C_1$, must contain infinitely many terms of the sequence. We select $\mathbf{x}_{k_1} \in C_1$. We repeat this process, constructing a nested sequence of non-empty compact sets $C_0 \supset C_1 \supset C_2 \supset \dots$ such that the diameter of $C_j$ tends to zero, and each $C_j$ contains infinitely many points of the sequence. We can then pick a subsequence $\{\mathbf{x}_{k_j}\}$ with $\mathbf{x}_{k_j} \in C_j$. This is a Cauchy sequence, and since $I^n$ is a complete metric space, the sequence converges to a limit $\mathbf{x}$. As each $C_j$ is closed, $\mathbf{x} \in C_j$ for all $j$, and thus $\mathbf{x} \in I^n$.
\end{proof}

\begin{lemma}[Uniform Continuity on a Compact Set]
    \label{lem:Uniform_Continuity}
    \uses{def:Continuous_Functions}
    \uses{lem:Bolzano_Weierstrass}
    Every function $f \in C(I^n)$ is uniformly continuous on the compact set $I^n$.
\end{lemma}

\begin{proof}
    Assume for contradiction that $f$ is not uniformly continuous. Then there exists an $\epsilon_0 > 0$ such that for every $k \in \mathbb{N}$, we can find points $\mathbf{x}_k, \mathbf{y}_k \in I^n$ with $\|\mathbf{x}_k - \mathbf{y}_k\|_\infty < 1/k$ but $|f(\mathbf{x}_k) - f(\mathbf{y}_k)| \ge \epsilon_0$.
    Since $I^n$ is compact, by Lemma \ref{lem:Bolzano_Weierstrass}, the sequence $\{\mathbf{x}_k\}$ has a convergent subsequence $\{\mathbf{x}_{k_j}\}$ with limit $\mathbf{x} \in I^n$. As $\|\mathbf{x}_{k_j} - \mathbf{y}_{k_j}\|_\infty < 1/k_j \to 0$, the corresponding subsequence $\{\mathbf{y}_{k_j}\}$ also converges to $\mathbf{x}$.
    By the continuity of $f$ at $\mathbf{x}$, we have $\lim_{j\to\infty} |f(\mathbf{x}_{k_j}) - f(\mathbf{y}_{k_j})| = |f(\mathbf{x}) - f(\mathbf{x})| = 0$. This contradicts that $|f(\mathbf{x}_{k_j}) - f(\mathbf{y}_{k_j})| \ge \epsilon_0$ for all $j$. Thus, $f$ must be uniformly continuous.
\end{proof}

\begin{lemma}[Piecewise-Constant Approximation]
    \label{lem:Piecewise_Constant_Approximation}
    The set of continuous, piecewise-constant functions on a compact interval $I = [a,b]$ is dense in $C(I)$. For any $f \in C(I)$ and any $\epsilon > 0$, there exists a continuous, piecewise-constant function $g$ such that $\|f - g\|_{\sup} < \epsilon$.
\end{lemma}

\begin{proof}
    Since $I$ is compact, $f$ is uniformly continuous. For a given $\epsilon > 0$, there exists a $\delta > 0$ such that $|x-y| < \delta$ implies $|f(x)-f(y)| < \epsilon/2$. Choose a partition of $I$, $a=x_0 < x_1 < \dots < x_m=b$, such that $\max_i(x_i - x_{i-1}) < \delta$. Define a piecewise-constant function $g_0(x) = f(x_i)$ for $x \in [x_i, x_{i+1})$. For any $x \in [x_i, x_{i+1})$, we have $|x - x_i| < \delta$, so $|f(x) - g_0(x)| = |f(x) - f(x_i)| < \epsilon/2$. Thus $\|f - g_0\|_{\sup} \le \epsilon/2$. To make the approximant continuous, we can modify it at the partition points to connect the constant segments with short linear ramps, ensuring the total deviation from $f$ remains less than $\epsilon$.
\end{proof}

\begin{lemma}[Linear Independence over Rationals]
    \label{lem:Linear_Independence_Over_Rationals}
    For any positive integer $m$, there exists a set of $m$ real numbers $\{\lambda_1, \dots, \lambda_m\}$ that is linearly independent over the field of rational numbers $\mathbb{Q}$.
\end{lemma}

\begin{proof}
    The real numbers $\mathbb{R}$ form a vector space over the field $\mathbb{Q}$. This vector space is of uncountable dimension. Any basis for this vector space must be uncountable. Therefore, we can always choose $m$ elements from this basis, which by definition are linearly independent over $\mathbb{Q}$. For a constructive example, one can choose numbers whose algebraic properties ensure independence, e.g., $\{1, \sqrt{p_1}, \sqrt{p_2}, \dots\}$ where $p_i$ are distinct primes.
\end{proof}

\begin{lemma}[Discretization into Addressable Grids]
    \label{lem:Addressable_Grids}
    \uses{def:Continuous_Functions}
    \uses{lem:Piecewise_Constant_Approximation}
    \uses{lem:Linear_Independence_Over_Rationals}
    For any set of continuous functions $\{\psi_q\}_{q=0}^{2n}$ in $C(I)$ and any $\epsilon > 0$, there exists a set of continuous, piecewise-constant functions $\{\phi_{q,p}\}_{q=0, p=1}^{2n, n}$ such that:
    \begin{enumerate}
        \item For each $q, p$, $\|\psi_q - \phi_{q,p}\| < \epsilon$.
        \item The sum $y_q(\mathbf{x}) = \lambda_q + \sum_{p=1}^n \phi_{q,p}(x_p)$, where $\{\lambda_q\}$ are linearly independent over $\mathbb{Q}$, defines a unique address for each rectangular block across all $2n+1$ grid systems.
    \end{enumerate}
\end{lemma}

\begin{proof}
    By Lemma \ref{lem:Piecewise_Constant_Approximation}, for each $\psi_q$, we can find a continuous, piecewise-constant function $\phi_{q,p}$ such that $\|\psi_q - \phi_{q,p}\| < \epsilon$. We choose the constant values taken by $\phi_{q,p}$ to be rational numbers.
    For each $q$, the sum $\sum_{p=1}^n \phi_{q,p}(x_p)$ is a rational value that is constant on each rectangular block of the $q$-th grid. By Lemma \ref{lem:Linear_Independence_Over_Rationals}, we can choose a set of $2n+1$ real numbers $\{\lambda_q\}_{q=0}^{2n}$ that are linearly independent over $\mathbb{Q}$.
    Define the address of a block as $y_q(\mathbf{x}) = \lambda_q + \sum_{p=1}^n \phi_{q,p}(x_p)$. Suppose two addresses are equal: $y_q(\mathbf{x}) = y_{q'}(\mathbf{x'})$. This implies $\lambda_q - \lambda_{q'} = \sum_{p=1}^n \phi_{q',p}(x'_p) - \sum_{p=1}^n \phi_{q,p}(x_p)$. The right side is a rational number. Due to the linear independence of $\{\lambda_q\}$, this equality can only hold if $q=q'$ (making $\lambda_q - \lambda_{q'} = 0$) and the rational parts are equal, which implies $\mathbf{x}$ and $\mathbf{x'}$ belong to the same block. Thus, each block across all grids has a unique address.
\end{proof}

\begin{lemma}[Grid Covering Property]
    \label{lem:Grid_Covering_Property}
    For any point $\mathbf{x} \in I^n$ and any $\delta > 0$, it is possible to construct $2n+1$ rectangular grid partitions of $I^n$ such that at least $n+1$ of the $2n+1$ unique grid blocks containing $\mathbf{x}$ (one from each partition) are fully contained within the open ball $B(\mathbf{x}, \delta)$.
\end{lemma}

\begin{proof}
    This is a non-trivial combinatorial geometric result central to Arnold's proof. The construction involves choosing the grid lines for the $2n+1$ partitions in a sufficiently independent or generic manner. A detailed proof is beyond the scope of this blueprint but is a foundational assumption for the subsequent steps.
\end{proof}

\begin{lemma}[Grid Sign Agreement Property]
    \label{lem:Grid_Sign_Agreement}
    \uses{lem:Uniform_Continuity}
    \uses{lem:Grid_Covering_Property}
    Let $f \in C(I^n)$ and $\mathbf{x} \in I^n$. For any $\epsilon > 0$ such that $|f(\mathbf{x})| > \epsilon$, one can construct $2n+1$ grid systems such that for at least $n+1$ indices $q$, the sign of $f$ is constant on the block containing $\mathbf{x}$.
\end{lemma}

\begin{proof}
    By Lemma \ref{lem:Uniform_Continuity}, $f$ is uniformly continuous. Given $\epsilon > 0$, we can find a $\delta > 0$ such that for all $\mathbf{y}$ in the ball $B(\mathbf{x}, \delta)$, we have $|f(\mathbf{y}) - f(\mathbf{x})| < \epsilon$. Since we assume $|f(\mathbf{x})| > \epsilon$, this implies that $\text{sgn}(f(\mathbf{y})) = \text{sgn}(f(\mathbf{x}))$ for all $\mathbf{y} \in B(\mathbf{x}, \delta)$.
    We then apply Lemma \ref{lem:Grid_Covering_Property}. We can construct $2n+1$ grids such that at least $n+1$ of the blocks containing $\mathbf{x}$ are fully contained within $B(\mathbf{x}, \delta)$. For these $n+1$ blocks, the sign of $f$ is therefore constant.
\end{proof}

\begin{lemma}[Single-Step Approximation]
    \label{lem:Single_Step_Approximation}
    \uses{def:Continuous_Functions}
    \uses{lem:Addressable_Grids}
    \uses{lem:Grid_Sign_Agreement}
    Let $f \in C(I^n)$ with $\|f\|=1$. For any $\delta \in (0, 1)$, there exist inner functions $\{\phi_{q,p}\}$ and a continuous outer function $\Phi$ such that $\left\|f - \sum_{q=0}^{2n} \Phi\left(\sum_{p=1}^n \phi_{q,p}(x_p)\right)\right\| \le 1 - \delta$.
\end{lemma}

\begin{proof}
    Choose a constant $C > 1/\delta$. Let $\epsilon_f = (C-1)/C$. For any point $\mathbf{x}$ with $|f(\mathbf{x})| \ge \epsilon_f$, we use Lemma \ref{lem:Grid_Sign_Agreement} to construct grids where at least $n+1$ blocks containing $\mathbf{x}$ have a constant sign for $f$. Let the unique addresses be defined as in Lemma \ref{lem:Addressable_Grids}. We define $\Phi$ at these addresses $a_{q,R}$:
    \[ \Phi(a_{q,R}) = \begin{cases} +1/C & \text{if } f(\mathbf{y}) > 0 \text{ for all } \mathbf{y} \in R \\ -1/C & \text{if } f(\mathbf{y}) < 0 \text{ for all } \mathbf{y} \in R \\ 0 & \text{otherwise} \end{cases} \]
    Extend $\Phi$ to a continuous function on $\mathbb{R}$ by linear interpolation. Let $F(\mathbf{x}) = \sum_{q=0}^{2n} \Phi\left(\sum_{p=1}^n \phi_{q,p}(x_p)\right)$. For any $\mathbf{x}$ where $f(\mathbf{x}) \ge \epsilon_f$, at least $n+1$ terms in the sum for $F(\mathbf{x})$ are $+1/C$ and at most $n$ are $-1/C$. So $F(\mathbf{x}) \ge (n+1)/C - n/C = 1/C$. The error is $|f(\mathbf{x}) - F(\mathbf{x})| \le 1 - 1/C < 1 - \delta$. A symmetric argument holds for $f(\mathbf{x}) \le -\epsilon_f$. The bound holds for all $\mathbf{x}$, thus $\|f - F\| \le 1 - \delta$.
\end{proof}

\begin{lemma}[Banach Space of Continuous Functions]
    \label{lem:Banach_Space_of_Continuous_Functions}
    \uses{def:Continuous_Functions}
    The space $C(I^n)$ of continuous functions on the compact set $I^n$ with the supremum norm is a complete metric space (a Banach space).
\end{lemma}

\begin{proof}
    Let $\{f_k\}$ be a Cauchy sequence in $C(I^n)$. For any $\epsilon > 0$, there exists an $N$ such that for $k,m > N$, $\|f_k - f_m\| < \epsilon$. For any fixed $\mathbf{x} \in I^n$, $|f_k(\mathbf{x}) - f_m(\mathbf{x})| < \epsilon$, so $\{f_k(\mathbf{x})\}$ is a Cauchy sequence of real numbers, which converges to a limit $f(\mathbf{x})$. This defines a function $f: I^n \to \mathbb{R}$.
    We must show convergence is uniform. Taking the limit as $m \to \infty$, we have $|f_k(\mathbf{x}) - f(\mathbf{x})| \le \epsilon$ for all $k>N$ and all $\mathbf{x} \in I^n$. This means $\|f_k - f\| \le \epsilon$ for $k>N$, so $f_k \to f$ uniformly.
    Finally, we show $f$ is continuous. Since $f$ is the uniform limit of a sequence of continuous functions, $f$ must be continuous. Therefore, $C(I^n)$ is complete.
\end{proof}

\begin{lemma}[Weierstrass M-Test]
    \label{lem:Weierstrass_M_Test}
    \uses{lem:Banach_Space_of_Continuous_Functions}
    Let $\{g_k\}$ be a sequence of functions in $C(X)$ where $X$ is a compact set. If there is a sequence of non-negative real numbers $\{M_k\}$ with $\|g_k\| \le M_k$ and $\sum M_k$ converges, then the series $\sum g_k$ converges uniformly to a continuous function.
\end{lemma}

\begin{proof}
    Let $S_N(\mathbf{x}) = \sum_{k=0}^N g_k(\mathbf{x})$ be the partial sum. For $M > N$, $\|S_M - S_N\| = \left\|\sum_{k=N+1}^M g_k\right\| \le \sum_{k=N+1}^M \|g_k\| \le \sum_{k=N+1}^M M_k$. Since $\sum M_k$ converges, its tail sum tends to 0. Thus, for any $\epsilon > 0$, there exists $N_0$ such that for $M > N > N_0$, $\sum_{k=N+1}^M M_k < \epsilon$. This shows that $\{S_N\}$ is a Cauchy sequence in $C(X)$. By Lemma \ref{lem:Banach_Space_of_Continuous_Functions}, $C(X)$ is a complete space, so the sequence of partial sums converges to a limit $S \in C(X)$.
\end{proof}

\begin{lemma}[Iterative Refinement to Exact Representation]
    \label{lem:Iterative_Refinement}
    \uses{lem:Single_Step_Approximation}
    \uses{lem:Weierstrass_M_Test}
    For any $f \in C(I^n)$, there exist inner functions $\{\phi_{q,p}\}$ and a continuous outer function $\Phi$ such that $f(\mathbf{x}) = \sum_{q=0}^{2n} \Phi\left(\sum_{p=1}^n \phi_{q,p}(x_p)\right)$.
\end{lemma}

\begin{proof}
    Let $f_0 = f$. Using Lemma \ref{lem:Single_Step_Approximation} with $\delta=1/2$, find $\Phi_0$ and $\{\phi_{q,p}\}$ such that $f_1 = f_0 - \sum_{q=0}^{2n} \Phi_0\left(\sum_{p=1}^n \phi_{q,p}(x_p)\right)$ has $\|f_1\| \le (1/2)\|f_0\|$. Repeat this process for $f_k$, defining $f_{k+1} = f_k - \sum_{q=0}^{2n} \Phi_k\left(\sum_{p=1}^n \phi_{q,p}(x_p)\right)$, yielding $\|f_{k+1}\| \le (1/2)\|f_k\|$. We get $\|f_k\| \le (1/2)^k \|f_0\| \to 0$.
    The construction of $\Phi_k$ in Lemma \ref{lem:Single_Step_Approximation} gives $\|\Phi_k\| \le \|f_k\|/C$ for some constant $C$. So, $\|\Phi_k\| \le M_k = (1/C)(1/2)^k \|f_0\|$. The series $\sum M_k$ converges. By Lemma \ref{lem:Weierstrass_M_Test}, the series $\Phi = \sum_{k=0}^{\infty} \Phi_k$ converges uniformly to a continuous function $\Phi$.
    The original function is a telescoping sum: $f = \sum_{k=0}^{\infty} (f_k - f_{k+1}) = \sum_{k=0}^{\infty} \sum_{q=0}^{2n} \Phi_k\left(\sum_{p=1}^n \phi_{q,p}(x_p)\right)$. Interchanging the sums gives $\sum_{q=0}^{2n} \left(\sum_{k=0}^{\infty} \Phi_k\right)\left(\sum_{p=1}^n \phi_{q,p}(x_p)\right) = \sum_{q=0}^{2n} \Phi\left(\sum_{p=1}^n \phi_{q,p}(x_p)\right)$.
\end{proof}

\begin{theorem}[Kolmogorov-Arnold Representation Theorem]
    \label{thm:Kolmogorov_Arnold}
    \uses{def:Continuous_Functions}
    \uses{lem:Iterative_Refinement}
    Every continuous function $f \in C(I^n)$ can be represented in the form:
    $$ f(x_1, \dots, x_n) = \sum_{q=0}^{2n} \Phi_q\left(\sum_{p=1}^n \phi_{q,p}(x_p)\right) $$
    where $\Phi_q$ are continuous functions of one variable, and $\phi_{q,p}$ are continuous functions of one variable.
\end{theorem}

\begin{proof}
    Lemma \ref{lem:Iterative_Refinement} proves that for any continuous function $f \in C(I^n)$, there exist a set of inner functions $\{\phi_{q,p}\}$ and a single continuous outer function $\Phi$ that provide an exact representation of the form $f(\mathbf{x}) = \sum_{q=0}^{2n} \Phi\left(\sum_{p=1}^n \phi_{q,p}(x_p)\right)$.
    This is a specific instance of the theorem's statement. By setting $\Phi_q = \Phi$ for all $q \in \{0, \dots, 2n\}$, we satisfy the conditions of the theorem. This establishes the existence of such a representation for any continuous function of $n$ variables.
\end{proof}

\chapter{Marden's Theorem}

\begin{definition}[Steiner Inellipse]
    \label{def:Steiner_Inellipse}
    For a triangle with vertices $z_1, z_2, z_3$ in the complex plane, the Steiner inellipse is the unique ellipse inscribed in the triangle and tangent to the sides at their midpoints.
\end{definition}

\begin{lemma}[Preservation of Geometric Properties under Affine Maps]
    \label{lem:Affine_Geometric_Preservation}
    An affine transformation $T(z) = az+b$ (with $a \neq 0$) preserves midpoints and tangency. That is, if $m$ is the midpoint of a segment $[p,q]$, then $T(m)$ is the midpoint of $[T(p), T(q)]$. If a line $L$ is tangent to a conic $C$ at point $p$, then the line $T(L)$ is tangent to the conic $T(C)$ at point $T(p)$.
\end{lemma}
\begin{proof}
    The midpoint of $[p,q]$ is $\frac{p+q}{2}$. Its image under $T$ is $a(\frac{p+q}{2}) + b = \frac{(ap+b)+(aq+b)}{2} = \frac{T(p)+T(q)}{2}$, which is the midpoint of the segment $[T(p), T(q)]$.
    Tangency is a geometric property defined by the intersection of a line and a conic having a single point with multiplicity two. Since affine transformations are defined by linear equations, they preserve the degree of polynomial equations defining curves and their intersections. Thus, a single point of intersection with multiplicity two is mapped to another such point.
\end{proof}

\begin{lemma}[Covariance of Ellipses and their Foci]
    \label{lem:Ellipse_Foci_Covariance}
    Let $E$ be an ellipse in the complex plane with foci $f_1$ and $f_2$. If $T(z) = az+b$ (with $a \neq 0$) is an affine transformation, then the image $T(E)$ is an ellipse whose foci are $T(f_1)$ and $T(f_2)$.
\end{lemma}
\begin{proof}
    An ellipse is a type of conic section, and affine transformations map conics to conics of the same type. Since an ellipse is a bounded conic, its image under an affine map must also be a bounded conic, which is an ellipse.
    The foci of an ellipse can be characterized by geometric properties that are preserved under affine maps. For instance, the foci are the singular points of the pencil of conics confocal with the given one. Affine maps preserve such pencils. A more direct argument is that any affine transformation can be decomposed into a rotation, a scaling, another rotation, and a translation. Each of these elementary transformations maps foci to foci. Thus, their composition does as well.
\end{proof}

\begin{lemma}[Covariance of the Steiner Inellipse]
    \label{lem:Steiner_Inellipse_Covariance}
    \uses{def:Steiner_Inellipse}
    \uses{lem:Affine_Geometric_Preservation}
    \uses{lem:Ellipse_Foci_Covariance}
    Let $\Delta$ be a triangle with vertices $z_1, z_2, z_3$ and let $S$ be its Steiner inellipse. For any affine transformation $T$, the image $T(S)$ is the Steiner inellipse of the transformed triangle $T(\Delta)$ with vertices $T(z_1), T(z_2), T(z_3)$.
\end{lemma}
\begin{proof}
    By lemma \ref{lem:Ellipse_Foci_Covariance}, $T(S)$ is an ellipse. By definition, $S$ is tangent to the midpoints of the sides of $\Delta$. By lemma \ref{lem:Affine_Geometric_Preservation}, an affine transformation preserves midpoints and tangency. Therefore, the ellipse $T(S)$ is tangent to the midpoints of the sides of the transformed triangle $T(\Delta)$. By the uniqueness of the Steiner inellipse (stated in definition \ref{def:Steiner_Inellipse}), $T(S)$ must be the Steiner inellipse of $T(\Delta)$.
\end{proof}

\begin{lemma}[Covariance of Roots and Derivative Roots]
    \label{lem:Root_Derivative_Covariance}
    Let $p(z)$ be a polynomial and let $T(w) = aw+b$ be an affine transformation with $a, b \in \mathbb{C}$ and $a \neq 0$. Let $q(w) = p(T(w))$. If $r$ is a root of $p(z)$ and $s$ is a root of its derivative $p'(z)$, then $T^{-1}(r)$ is a root of $q(w)$ and $T^{-1}(s)$ is a root of its derivative $q'(w)$.
\end{lemma}
\begin{proof}
    If $r$ is a root of $p(z)$, then $p(r)=0$. The roots of $q(w)$ are the values of $w$ such that $p(T(w))=0$, which means $T(w)=r$, or $w=T^{-1}(r)$.
    By the chain rule, the derivative of $q(w)$ is $q'(w) = a \cdot p'(T(w))$. The roots of $q'(w)$ are the values of $w$ for which $p'(T(w))=0$. If $s$ is a root of $p'(z)$, then we must have $T(w)=s$, which means $w=T^{-1}(s)$.
\end{proof}

\begin{lemma}[Normalization of a Triangle]
    \label{lem:Normalization}
    \uses{def:Steiner_Inellipse}
    For any non-collinear set of points $\{z_1, z_2, z_3\}$, there exists an affine transformation that maps them to points $\{w_1, w_2, w_3\}$ forming a triangle whose Steiner inellipse is centered at the origin with its foci on the real axis at $\pm c$ for some $c \ge 0$.
\end{lemma}
\begin{proof}
    Let $\Delta$ be the original triangle. Its Steiner inellipse $S$ has a center $m$ and foci $f_1, f_2$. First, apply a translation $T_1(z) = z - m$ to center the ellipse at the origin. Let the foci of the new ellipse be $f'_1, f'_2$. Second, apply a rotation $T_2(z) = z e^{-i\arg(f'_1)}$ to align the foci with the real axis. The composition $T = T_2 \circ T_1$ is an affine transformation that maps the triangle to the desired standard position. The foci will be at $\pm c$ for some $c = |f'_1| \ge 0$.
\end{proof}

\begin{lemma}[Canonical Vertices for a Standard Triangle]
    \label{lem:Vertex_Representation_Standard}
    Let $\Delta$ be a triangle whose Steiner inellipse is centered at the origin and has foci $\pm c$ on the real axis. Let the semi-major axis be $a$ and semi-minor axis be $b$, so $c^2=a^2-b^2$. The vertices of such a triangle can be represented as $\{2a, -a+ib\sqrt{3}, -a-ib\sqrt{3}\}$ up to a rotation by $\pi$.
\end{lemma}
\begin{proof}
    This is a known parameterization. It can be derived by applying the affine transformation $T(x+iy) = ax+iby$ to an equilateral triangle with vertices $\{2, -1+i\sqrt{3}, -1-i\sqrt{3}\}$. The Steiner inellipse of this equilateral triangle is the unit circle, which is transformed by $T$ into an ellipse with semi-axes $a$ and $b$ aligned with the coordinate axes. The vertices of the transformed triangle are $T(2) = 2a$, $T(-1+i\sqrt{3}) = -a+ib\sqrt{3}$, and $T(-1-i\sqrt{3}) = -a-ib\sqrt{3}$. The centroid of these vertices is $\frac{1}{3}(2a -a+ib\sqrt{3} -a-ib\sqrt{3}) = 0$.
\end{proof}

\begin{lemma}[Sum of Squared Vertices for a Standard Triangle]
    \label{lem:Sum_Squared_Vertices_Standard}
    \uses{lem:Vertex_Representation_Standard}
    For a triangle in the standard position described in lemma \ref{lem:Vertex_Representation_Standard}, with Steiner inellipse foci at $\pm c$, the sum of the squares of its vertices is $6c^2$.
\end{lemma}
\begin{proof}
    Using the vertex representation from lemma \ref{lem:Vertex_Representation_Standard}, $w_1=2a, w_2=-a+ib\sqrt{3}, w_3=-a-ib\sqrt{3}$.
    \begin{align*}
        w_1^2 &= (2a)^2 = 4a^2 \\
        w_2^2 &= (-a+ib\sqrt{3})^2 = a^2 - 3b^2 - 2iab\sqrt{3} \\
        w_3^2 &= (-a-ib\sqrt{3})^2 = a^2 - 3b^2 + 2iab\sqrt{3}
    \end{align*}
    Summing these gives:
    $\sum_{j=1}^3 w_j^2 = 4a^2 + (a^2 - 3b^2) + (a^2 - 3b^2) = 6a^2 - 6b^2 = 6(a^2-b^2)$.
    Since the foci are at $\pm c$ where $c^2=a^2-b^2$, we have $\sum_{j=1}^3 w_j^2 = 6c^2$.
\end{proof}

\begin{lemma}[Roots of the Derivative for Centered Polynomials]
    \label{lem:Vieta_Centered}
    Let $p(z)$ be a monic cubic polynomial with roots $z_1, z_2, z_3$ satisfying $\sum_{j=1}^3 z_j = 0$. The roots of the derivative $p'(z)$ are given by $\pm\sqrt{-\frac{1}{3}\sum_{j<k} z_j z_k}$.
\end{lemma}
\begin{proof}
    Given $\sum z_j = 0$, the polynomial is $p(z) = z^3 + (z_1z_2+z_2z_3+z_3z_1)z - z_1z_2z_3$.
    The derivative is $p'(z) = 3z^2 + (z_1z_2+z_2z_3+z_3z_1)$. The roots of $p'(z)$ are the values of $z$ such that $3z^2 = -(z_1z_2+z_2z_3+z_3z_1)$.
    Thus, $z = \pm\sqrt{-\frac{1}{3}\sum_{j<k} z_j z_k}$.
\end{proof}

\begin{lemma}[Foci as Derivative Roots for Standard Triangle]
    \label{lem:Foci_Standard_Case}
    \uses{lem:Sum_Squared_Vertices_Standard}
    \uses{lem:Vieta_Centered}
    For a triangle in the standard position (centroid at origin, Steiner inellipse foci at $\pm c$), the roots of the derivative of the polynomial whose roots are the vertices are precisely the foci of the Steiner inellipse.
\end{lemma}
\begin{proof}
    Let the vertices be $w_1, w_2, w_3$. Since the centroid is at the origin, $\sum w_j=0$. By lemma \ref{lem:Vieta_Centered}, the roots of the derivative are $\pm\sqrt{-\frac{1}{3}\sum_{j<k} w_j w_k}$.
    From the identity $(\sum w_j)^2 = \sum w_j^2 + 2\sum_{j<k} w_j w_k$, and since $\sum w_j=0$, we have $\sum_{j<k} w_j w_k = -\frac{1}{2}\sum w_j^2$.
    The roots of the derivative are therefore $\pm\sqrt{-\frac{1}{3}(-\frac{1}{2}\sum w_j^2)} = \pm\sqrt{\frac{1}{6}\sum w_j^2}$.
    By lemma \ref{lem:Sum_Squared_Vertices_Standard}, for a triangle in standard position, $\sum w_j^2 = 6c^2$.
    Substituting this result, the roots of the derivative are $\pm\sqrt{\frac{1}{6}(6c^2)} = \pm\sqrt{c^2} = \pm c$. These are precisely the foci of the Steiner inellipse.
\end{proof}

\begin{theorem}[Marden's Theorem]
    \label{thm:Marden_Theorem}
    \uses{def:Steiner_Inellipse}
    \uses{lem:Steiner_Inellipse_Covariance}
    \uses{lem:Root_Derivative_Covariance}
    \uses{lem:Normalization}
    \uses{lem:Foci_Standard_Case}
    Suppose the zeroes $z_1, z_2, z_3$ of a third-degree polynomial $p(z)$ are non-collinear complex numbers. The foci of the Steiner inellipse of the triangle with vertices $z_1, z_2, z_3$ are the zeroes of the derivative $p'(z)$.
\end{theorem}
\begin{proof}
    Let $\Delta_z$ be the triangle with vertices $z_1, z_2, z_3$. By lemma \ref{lem:Normalization}, there exists an affine transformation $T$ that maps the vertices of a triangle $\Delta_w$ in standard position to the vertices of $\Delta_z$. Let $w_1, w_2, w_3$ be the vertices of $\Delta_w$ and $f_w, g_w$ be the foci of its Steiner inellipse. Then $z_j = T(w_j)$.
    Let $q(w)$ be the polynomial with roots $w_1, w_2, w_3$. By lemma \ref{lem:Foci_Standard_Case}, the roots of $q'(w)$ are the foci $f_w, g_w$.
    Let $p(z)$ be the polynomial with roots $z_j$. Then $p(z)$ is proportional to $q(T^{-1}(z))$.
    By lemma \ref{lem:Root_Derivative_Covariance}, the roots of $p'(z)$ are the images of the roots of $q'(w)$ under $T$. So the roots of $p'(z)$ are $T(f_w)$ and $T(g_w)$.
    By lemma \ref{lem:Steiner_Inellipse_Covariance}, the Steiner inellipse of $\Delta_z$ is the image under $T$ of the Steiner inellipse of $\Delta_w$. The foci of the Steiner inellipse of $\Delta_z$ are the images under $T$ of the foci of the Steiner inellipse of $\Delta_w$, which are $T(f_w)$ and $T(g_w)$.
    Therefore, the roots of $p'(z)$ are the foci of the Steiner inellipse of $\Delta_z$.
\end{proof}

\chapter{Menelaus's Theorem}

\begin{definition}[Signed Length of a Segment]
    \label{def:Signed_Length}
    For any two distinct points A and B on an oriented line, the signed length $\overline{AB}$ is the distance between A and B, taken as positive if the direction from A to B is the same as the orientation of the line, and negative otherwise. For any three collinear points A, B, and C, this implies $\overline{AB} = -\overline{BA}$ and $\overline{AC} = \overline{AB} + \overline{BC}$. A ratio $\frac{\overline{AP}}{\overline{PB}}$ is positive if P is between A and B, and negative otherwise.
\end{definition}

\begin{definition}[Transversal of a Triangle]
    \label{def:Transversal}
    A transversal of a triangle $\triangle ABC$ is a line that intersects the lines containing the sides of the triangle, at points distinct from the vertices A, B, and C.
\end{definition}

\begin{lemma}[AA Similarity Criterion]
    \label{lem:AA_Similarity}
    If two angles of one triangle are congruent to two angles of another triangle, then the two triangles are similar. This implies that the ratio of corresponding side lengths is constant.
\end{lemma}

\begin{lemma}[Segment Division and Plane Separation]
    \label{lem:Plane_Separation_Consequence}
    Let A and B be distinct points in a plane, and let a transversal line $l$ intersect the line AB at a point F. Then F lies between A and B if and only if A and B lie on opposite sides of the line $l$. Consequently, if F is outside the segment AB, then A and B lie on the same side of $l$.
\end{lemma}

\begin{proof}
    This is a consequence of the Plane Separation Postulate, which states that any line in a plane divides the plane into two disjoint convex sets (half-planes). If two points lie in different half-planes, the segment connecting them must intersect the line. If F is between A and B, the segment AB contains F, which is on line $l$. Since F is the only intersection point of line AB and line $l$, A and B cannot both be in the same half-plane, as that would imply the segment AB is entirely contained within that half-plane, a contradiction. Thus, A and B are on opposite sides. Conversely, if A and B are on opposite sides of $l$, the segment AB must intersect $l$ at some point, which must be F. Therefore, F is on the segment AB.
\end{proof}

\begin{lemma}[Segment Ratio by Signed Perpendiculars]
    \label{lem:Segment_Ratio_by_Signed_Perpendiculars}
    \uses{def:Signed_Length}
    \uses{lem:AA_Similarity}
    \uses{lem:Plane_Separation_Consequence}
    Let A and B be two distinct points, and let a line $l$ intersect the line AB at a point F. Let $h_A$ and $h_B$ be the signed distances from A and B to the line $l$, where distances on opposite sides of $l$ have opposite signs. Then, the ratio of the signed lengths is given by:
    $$ \frac{\overline{AF}}{\overline{FB}} = -\frac{h_A}{h_B} $$
\end{lemma}

\begin{proof}
    Let $P_A$ and $P_B$ be the feet of the perpendiculars from A and B to line $l$. The triangles $\triangle AFP_A$ and $\triangle BFP_B$ are right-angled triangles. The angles $\angle AFP_A$ and $\angle BFP_B$ are vertically opposite, so they are equal. By lemma \ref{lem:AA_Similarity}, $\triangle AFP_A \sim \triangle BFP_B$. Therefore, the ratio of the absolute lengths of corresponding sides is equal:
    $$ \frac{|\overline{AF}|}{|\overline{FB}|} = \frac{|h_A|}{|h_B|} $$
    Now we analyze the signs.
    \begin{itemize}
        \item Case 1: F is between A and B. By definition \ref{def:Signed_Length}, the ratio $\frac{\overline{AF}}{\overline{FB}}$ is positive. By lemma \ref{lem:Plane_Separation_Consequence}, A and B are on opposite sides of $l$, so their signed distances $h_A$ and $h_B$ have opposite signs, which means the ratio $\frac{h_A}{h_B}$ is negative.
        \item Case 2: F is outside the segment AB. By definition \ref{def:Signed_Length}, the ratio $\frac{\overline{AF}}{\overline{FB}}$ is negative. By lemma \ref{lem:Plane_Separation_Consequence}, A and B are on the same side of $l$, so their signed distances $h_A$ and $h_B$ have the same sign, which means the ratio $\frac{h_A}{h_B}$ is positive.
    \end{itemize}
    In both cases, the ratio of signed lengths $\frac{\overline{AF}}{\overline{FB}}$ has the opposite sign to the ratio of signed distances $\frac{h_A}{h_B}$. Since their magnitudes are equal, we conclude:
    $$ \frac{\overline{AF}}{\overline{FB}} = -\frac{h_A}{h_B} $$
\end{proof}

\begin{theorem}[Menelaus's Theorem]
    \label{thm:Menelaus}
    \uses{def:Signed_Length}
    \uses{def:Transversal}
    \uses{lem:Segment_Ratio_by_Signed_Perpendiculars}
    Let $\triangle ABC$ be a triangle, and let a line be a transversal that intersects the lines containing the sides AB, BC, and CA at points F, D, and E, respectively. Then,
    $$ \frac{\overline{AF}}{\overline{FB}} \cdot \frac{\overline{BD}}{\overline{DC}} \cdot \frac{\overline{CE}}{\overline{EA}} = -1 $$
\end{theorem}

\begin{proof}
    \uses{lem:Segment_Ratio_by_Signed_Perpendiculars}
    Let the transversal line be denoted by $l$. Let $h_A$, $h_B$, and $h_C$ be the signed distances from the vertices A, B, and C to the line $l$.
    By applying lemma \ref{lem:Segment_Ratio_by_Signed_Perpendiculars} to each side of the triangle, we obtain the following three relations:
    \begin{align*}
        \frac{\overline{AF}}{\overline{FB}} &= -\frac{h_A}{h_B} \\
        \frac{\overline{BD}}{\overline{DC}} &= -\frac{h_B}{h_C} \\
        \frac{\overline{CE}}{\overline{EA}} &= -\frac{h_C}{h_A}
    \end{align*}
    Multiplying these three equations together yields:
    $$ \frac{\overline{AF}}{\overline{FB}} \cdot \frac{\overline{BD}}{\overline{DC}} \cdot \frac{\overline{CE}}{\overline{EA}} = \left(-\frac{h_A}{h_B}\right) \cdot \left(-\frac{h_B}{h_C}\right) \cdot \left(-\frac{h_C}{h_A}\right) = - \frac{h_A h_B h_C}{h_B h_C h_A} = -1 $$
    This completes the proof.
\end{proof}

\begin{lemma}[Uniqueness of a Point Dividing a Segment in a Given Ratio]
    \label{lem:Uniqueness_of_Division_Point}
    \uses{def:Signed_Length}
    Let A and B be two distinct points on a line. For any real number $k \neq -1$, there exists a unique point F on the line AB such that $\frac{\overline{AF}}{\overline{FB}} = k$.
\end{lemma}

\begin{proof}
    Let the line be represented by the real number line, with the coordinate of A being $a$ and the coordinate of B being $b$. Let the coordinate of F be $x$. The signed lengths are $\overline{AF} = x-a$ and $\overline{FB} = b-x$. The condition is:
    $$ \frac{x-a}{b-x} = k $$
    We solve for $x$:
    $$ x - a = k(b - x) \implies x - a = kb - kx \implies x(1+k) = a + kb $$
    Since $k \neq -1$, the term $(1+k)$ is non-zero, and we can divide by it to find a unique value for $x$:
    $$ x = \frac{a+kb}{1+k} $$
    This shows that for any $k \neq -1$, there is a unique point F that satisfies the condition.
\end{proof}

\begin{lemma}[Non-parallelism Condition for Converse Menelaus]
    \label{lem:Non_Parallelism}
    \uses{def:Signed_Length}
    Let $\triangle ABC$ be a triangle and let D, E, F be points on lines BC, CA, AB respectively, satisfying the Menelaus condition $\frac{\overline{AF}}{\overline{FB}} \cdot \frac{\overline{BD}}{\overline{DC}} \cdot \frac{\overline{CE}}{\overline{EA}} = -1$. Then the line DE cannot be parallel to the line AB.
\end{lemma}

\begin{proof}
    \uses{def:Signed_Length}
    Assume for contradiction that the line DE is parallel to the line AB. By the property of similar triangles ($\triangle CDE \sim \triangle CBA$), we have $\frac{\overline{CE}}{\overline{EA}} = \frac{\overline{CD}}{\overline{DB}}$.
    This can be rewritten as $\frac{\overline{CE}}{\overline{EA}} = \frac{-\overline{DC}}{-\overline{BD}} = \frac{\overline{DC}}{\overline{BD}}$.
    Therefore, $\frac{\overline{BD}}{\overline{DC}} \cdot \frac{\overline{CE}}{\overline{EA}} = 1$.
    Substituting this into the given Menelaus condition gives:
    $$ \frac{\overline{AF}}{\overline{FB}} \cdot \left( \frac{\overline{BD}}{\overline{DC}} \cdot \frac{\overline{CE}}{\overline{EA}} \right) = -1 \implies \frac{\overline{AF}}{\overline{FB}} \cdot (1) = -1 $$
    This implies $\frac{\overline{AF}}{\overline{FB}} = -1$, which means $\overline{AF} = -\overline{FB}$. By the property of signed lengths, $\overline{AF} = \overline{AB} + \overline{BF} = \overline{AB} - \overline{FB}$.
    Substituting $\overline{AF} = -\overline{FB}$ gives $\overline{AB} - \overline{FB} = -\overline{FB}$, which simplifies to $\overline{AB} = 0$. This implies that A and B are the same point, which contradicts the definition of a triangle.
    Thus, our initial assumption that DE is parallel to AB must be false.
\end{proof}

\begin{theorem}[Converse of Menelaus's Theorem]
    \label{thm:Converse_of_Menelaus}
    \uses{def:Signed_Length}
    \uses{thm:Menelaus}
    \uses{lem:Uniqueness_of_Division_Point}
    \uses{lem:Non_Parallelism}
    Let $\triangle ABC$ be a triangle, and let D, E, and F be points on the lines containing the sides BC, CA, and AB, respectively. If
    $$ \frac{\overline{AF}}{\overline{FB}} \cdot \frac{\overline{BD}}{\overline{DC}} \cdot \frac{\overline{CE}}{\overline{EA}} = -1 $$
    then the points D, E, and F are collinear.
\end{theorem}

\begin{proof}
    By lemma \ref{lem:Non_Parallelism}, the line DE is not parallel to the line AB. Therefore, the line passing through D and E must intersect the line AB at exactly one point. Let's call this point F'.
    Since D, E, and F' are collinear, by theorem \ref{thm:Menelaus}, they must satisfy the equation:
    $$ \frac{\overline{AF'}}{\overline{F'B}} \cdot \frac{\overline{BD}}{\overline{DC}} \cdot \frac{\overline{CE}}{\overline{EA}} = -1 $$
    We are given that the point F satisfies:
    $$ \frac{\overline{AF}}{\overline{FB}} \cdot \frac{\overline{BD}}{\overline{DC}} \cdot \frac{\overline{CE}}{\overline{EA}} = -1 $$
    Since D and E are not vertices, the terms $\frac{\overline{BD}}{\overline{DC}}$ and $\frac{\overline{CE}}{\overline{EA}}$ are well-defined, finite, and non-zero. Comparing the two equations gives:
    $$ \frac{\overline{AF'}}{\overline{F'B}} = \frac{\overline{AF}}{\overline{FB}} $$
    Let this ratio be $k$. Lemma \ref{lem:Non_Parallelism} ensures that $k \neq -1$.
    By lemma \ref{lem:Uniqueness_of_Division_Point}, there is only one point on the line AB that divides the segment AB in a given ratio $k$. Therefore, F and F' must be the same point.
    Since F' is on the line DE by construction, F must also be on the line DE. Thus, the points D, E, and F are collinear.
\end{proof} 

\chapter{Pick's Theorem}

\begin{definition}[Simple Polygon]
    \label{def:Simple_Polygon}
    A simple polygon is a polygon in the plane whose boundary does not cross itself. The vertices of the polygon are points on an integer lattice, i.e., points with integer coordinates.
\end{definition}

\begin{lemma}[Additivity of Area]
    \label{lem:Additivity_of_Area}
    \uses{def:Simple_Polygon}
    Let a simple polygon $P$ be partitioned into a finite set of non-overlapping simple polygons $\{P_j\}$. Then the area of $P$ is the sum of the areas of the polygons $P_j$, i.e., $A(P) = \sum_{j} A(P_j)$.
\end{lemma}

\begin{proof}
    This is a fundamental property of the geometric area of planar figures. If a region is the disjoint union of a finite number of sub-regions (up to their boundaries), its total area is the sum of the areas of the sub-regions.
\end{proof}

\begin{lemma}[Area of a Rectangle with Integer Vertices]
    \label{lem:Area_of_a_Rectangle}
    \uses{def:Simple_Polygon}
    For a simple rectangle with vertices at integer coordinates and sides parallel to the coordinate axes, the area $A$ is given by $A = i + \frac{b}{2} - 1$, where $i$ is the number of interior integer points and $b$ is the number of integer points on the boundary.
\end{lemma}

\begin{proof}
    Let the rectangle have vertices at $(x_1,y_1), (x_2,y_1), (x_1,y_2), (x_2,y_2)$ for integers $x_1<x_2$ and $y_1<y_2$. Let $w = x_2-x_1$ and $h=y_2-y_1$. The area of the rectangle is $A = wh$.
    The number of interior points is $i = (w-1)(h-1)$.
    The number of boundary points is $b = 2w + 2h$.
    We verify the formula:
    \begin{align*}
        i + \frac{b}{2} - 1 &= (w-1)(h-1) + \frac{2w+2h}{2} - 1 \\
        &= (wh - w - h + 1) + (w+h) - 1 \\
        &= wh = A.
    \end{align*}
    Thus, the formula holds.
\end{proof}

\begin{lemma}[Central Symmetry of Rectangle Triangles]
    \label{lem:Central_Symmetry_of_Rectangle_Triangles}
    Let $R$ be an axis-aligned rectangle. Let a diagonal $D$ divide $R$ into two congruent right-angled triangles $T_1$ and $T_2$. The central symmetry transformation with respect to the rectangle's center is a bijection from the interior of $T_1$ to the interior of $T_2$.
\end{lemma}

\begin{proof}
    Let the vertices of the rectangle be $(x_1, y_1), (x_2, y_1), (x_1, y_2), (x_2, y_2)$. The center of the rectangle is $C = (\frac{x_1+x_2}{2}, \frac{y_1+y_2}{2})$. The central symmetry transformation with respect to $C$ is $S(p) = 2C - p$.
    This transformation maps the rectangle $R$ to itself. Specifically, it maps vertices to opposite vertices. For instance, $S((x_1,y_1)) = (x_2,y_2)$.
    The diagonal $D$ is mapped to itself. A point $p$ is in the interior of $T_1$ if and only if it is in the interior of $R$ and on one side of the diagonal $D$. The map $S$ preserves the interior of $R$ and maps points from one side of $D$ to the other. Thus, $S$ maps the interior of $T_1$ to the interior of $T_2$.
    Since $S(S(p)) = p$, the transformation is its own inverse, and is therefore a bijection.
\end{proof}

\begin{lemma}[Symmetry of Lattice Points in a Diagonal-Cut Rectangle]
    \label{lem:Symmetry_of_Lattice_Points}
    \uses{lem:Central_Symmetry_of_Rectangle_Triangles}
    Let $R$ be a rectangle with integer vertices and sides parallel to the coordinate axes. A diagonal of $R$ divides it into two congruent right-angled triangles, $T_1$ and $T_2$. The number of interior integer points in $T_1$ is equal to the number of interior integer points in $T_2$.
\end{lemma}

\begin{proof}
    Let the vertices of the rectangle be $(x_1, y_1), (x_2, y_1), (x_1, y_2), (x_2, y_2)$, where $x_1, x_2, y_1, y_2$ are integers. The central symmetry transformation is $S(x,y) = (x_1+x_2-x, y_1+y_2-y)$.
    If a point $(x,y)$ has integer coordinates, then $S(x,y)$ also has integer coordinates.
    By lemma \ref{lem:Central_Symmetry_of_Rectangle_Triangles}, $S$ is a bijection between the interiors of $T_1$ and $T_2$.
    Since $S$ preserves integer coordinates, it establishes a one-to-one correspondence between the set of interior integer points of $T_1$ and the set of interior integer points of $T_2$. Consequently, their counts must be equal.
\end{proof}

\begin{lemma}[Area of a Right-Angled Triangle]
    \label{lem:Area_of_a_Right_Angled_Triangle}
    \uses{lem:Area_of_a_Rectangle}
    \uses{lem:Symmetry_of_Lattice_Points}
    For a right-angled triangle with integer vertices and legs parallel to the coordinate axes, the area $A$ is given by $A = i + \frac{b}{2} - 1$.
\end{lemma}

\begin{proof}
    Let the right-angled triangle $T$ be formed by cutting a rectangle $R$ along its diagonal. Let $A_R, i_R, b_R$ be the area, interior points, and boundary points for $R$, and $A_T, i_T, b_T$ be for $T$.
    By lemma \ref{lem:Area_of_a_Rectangle}, $A_R = i_R + \frac{b_R}{2} - 1$.
    Let $k$ be the number of integer points on the diagonal. Let the rectangle have width $w$ and height $h$. The area is $A_T = \frac{1}{2} A_R$.
    By lemma \ref{lem:Symmetry_of_Lattice_Points}, the number of interior points in the two triangles formed by the diagonal are equal. Thus, the interior points of the rectangle consist of the interior points of the two triangles plus the points on the diagonal (excluding vertices): $i_R = 2i_T + (k-2)$.
    The boundary points of the triangle are $b_T = w+h+k-1$. From the proof of lemma \ref{lem:Area_of_a_Rectangle}, $b_R=2w+2h$.
    First, we express $A_T$ in terms of $i_T$ and the geometric parameters:
    \begin{align*}
        A_T &= \frac{1}{2}A_R = \frac{1}{2} \left(i_R + \frac{b_R}{2} - 1\right) \\
            &= \frac{1}{2} \left( (2i_T + k - 2) + \frac{2w+2h}{2} - 1 \right) \\
            &= i_T + \frac{k-2}{2} + \frac{w+h}{2} - \frac{1}{2} \\
            &= i_T + \frac{w+h+k-3}{2}.
    \end{align*}
    Now, we evaluate Pick's formula for the triangle:
    \begin{align*}
        i_T + \frac{b_T}{2} - 1 &= i_T + \frac{w+h+k-1}{2} - 1 \\
                               &= i_T + \frac{w+h+k-1-2}{2} \\
                               &= i_T + \frac{w+h+k-3}{2}.
    \end{align*}
    The expressions are identical, so $A_T = i_T + \frac{b_T}{2} - 1$.
\end{proof}

\begin{lemma}[Additivity of Pick's Formula Expression]
    \label{lem:Additivity_of_Picks_Formula}
    \uses{def:Simple_Polygon}
    Let a simple polygon $P$ be formed by joining two simple polygons $P_1$ and $P_2$ along a common boundary edge. Let $F(S) = i_S + \frac{b_S}{2} - 1$ for any polygon $S$. Then $F(P) = F(P_1) + F(P_2)$.
\end{lemma}

\begin{proof}
    Let $i_1, b_1$ be the counts for $P_1$, and $i_2, b_2$ for $P_2$. Let $i, b$ be for $P$.
    Let the shared boundary contain $k$ integer points (including endpoints).
    The interior points of $P$ are $i = i_1 + i_2 + (k-2)$.
    The boundary points of $P$ are $b = b_1 + b_2 - 2(k-1) - 2 = b_1 + b_2 - 2k + 2$.
    Now, we compute $F(P_1) + F(P_2)$ and substitute these relations:
    \begin{align*}
        F(P_1) + F(P_2) &= \left(i_1 + \frac{b_1}{2} - 1\right) + \left(i_2 + \frac{b_2}{2} - 1\right) \\
        &= (i_1 + i_2) + \frac{b_1+b_2}{2} - 2 \\
        &= (i-k+2) + \frac{b+2k-2}{2} - 2 \\
        &= i-k+2 + \frac{b}{2} + k - 1 - 2 \\
        &= i + \frac{b}{2} - 1 = F(P).
    \end{align*}
    Thus, the function $F$ is additive.
\end{proof}

\begin{lemma}[Decomposition of a Triangle into a Rectangle and Right Triangles]
    \label{lem:Triangle_in_Rectangle_Decomposition}
    Any triangle $T$ with integer vertices can be enclosed in an axis-aligned rectangle $R$ with integer vertices, such that the region $R \setminus T$ is composed of at most three right-angled triangles whose vertices are also integer points and whose legs are parallel to the coordinate axes.
\end{lemma}

\begin{proof}
    Let the vertices of the triangle $T$ be $(x_1, y_1), (x_2, y_2), (x_3, y_3)$.
    Define the bounding rectangle $R$ by the minimum and maximum coordinates: $x_{\min} = \min(x_i)$, $x_{\max} = \max(x_i)$, $y_{\min} = \min(y_i)$, $y_{\max} = \max(y_i)$. The vertices of $R$ are integer points.
    The region $R \setminus T$ is composed of corner regions. Each of these regions is either empty or a right-angled triangle with integer vertices and legs parallel to the coordinate axes. There can be at most three such triangles.
\end{proof}

\begin{lemma}[Area of a General Triangle]
    \label{lem:Area_of_a_General_Triangle}
    \uses{lem:Area_of_a_Rectangle}
    \uses{lem:Area_of_a_Right_Angled_Triangle}
    \uses{lem:Additivity_of_Picks_Formula}
    \uses{lem:Additivity_of_Area}
    \uses{lem:Triangle_in_Rectangle_Decomposition}
    For any triangle with integer vertices, the area $A$ is given by $A = i + \frac{b}{2} - 1$.
\end{lemma}

\begin{proof}
    Let $T$ be a general triangle. By lemma \ref{lem:Triangle_in_Rectangle_Decomposition}, $T$ can be embedded in a rectangle $R$ such that $R$ is the disjoint union of $T$ and a set of at most three right-angled triangles, $\{T_j\}$.
    Let $F(P) = i_P + \frac{b_P}{2} - 1$. By lemma \ref{lem:Additivity_of_Picks_Formula}, the function $F$ is additive, so $F(R) = F(T) + \sum F(T_j)$.
    From lemma \ref{lem:Area_of_a_Rectangle}, we know $A(R) = F(R)$.
    From lemma \ref{lem:Area_of_a_Right_Angled_Triangle}, we know $A(T_j) = F(T_j)$ for each corner triangle $T_j$.
    Substituting these identities into the additivity equation for $F$ gives: $A(R) = F(T) + \sum A(T_j)$.
    By lemma \ref{lem:Additivity_of_Area}, we have $A(R) = A(T) + \sum A(T_j)$.
    Comparing the two expressions for $A(R)$, we conclude that $A(T) = F(T)$.
\end{proof}

\begin{lemma}[Existence of an Internal Diagonal]
    \label{lem:Existence_of_Internal_Diagonal}
    \uses{def:Simple_Polygon}
    Every simple polygon with $n \ge 4$ vertices has at least one internal diagonal.
\end{lemma}

\begin{proof}
    Let $P$ be a simple polygon. A vertex $v$ is convex if its interior angle is less than $\pi$. Every polygon has at least three convex vertices. Let $v$ be such a convex vertex, and let $u$ and $w$ be its adjacent vertices. Consider the triangle $\triangle uvw$.
    If the segment $uw$ is an internal diagonal, we are done. This happens if the interior of $uw$ is in the interior of $P$.
    If $uw$ is not an internal diagonal, there must be one or more vertices of $P$ in the interior of $\triangle uvw$. Let $z$ be such a vertex that is farthest from the line segment $uw$. The line segment $vz$ must lie entirely inside the polygon. If it did not, it would have to cross an edge of $P$, which would imply another vertex is inside $\triangle uvw$ and farther from $uw$ than $z$, a contradiction. Therefore, $vz$ is an internal diagonal of $P$.
\end{proof}

\begin{lemma}[Polygon Triangulation Theorem]
    \label{lem:Polygon_Triangulation}
    \uses{def:Simple_Polygon}
    \uses{lem:Existence_of_Internal_Diagonal}
    Every simple polygon with $n \ge 3$ vertices can be partitioned into $n-2$ non-overlapping triangles whose vertices are the vertices of the original polygon.
\end{lemma}

\begin{proof}
    We prove this by induction on the number of vertices $n$.
    Base Case ($n=3$): The polygon is a triangle. It is already triangulated into $3-2=1$ triangle.
    Inductive Step: Assume any simple polygon with $k < n$ vertices can be triangulated into $k-2$ triangles. Let $P$ be a polygon with $n > 3$ vertices. By lemma \ref{lem:Existence_of_Internal_Diagonal}, there exists an internal diagonal. This diagonal splits $P$ into two smaller simple polygons, $P_1$ and $P_2$, with $n_1$ and $n_2$ vertices, respectively. We have $n_1, n_2 < n$ and $n_1+n_2 = n+2$.
    By the induction hypothesis, $P_1$ can be triangulated into $n_1-2$ triangles and $P_2$ into $n_2-2$ triangles. The union of these triangulations forms a triangulation of $P$.
    The total number of triangles is $(n_1-2) + (n_2-2) = n_1+n_2-4 = (n+2)-4 = n-2$.
\end{proof}

\begin{theorem}[Pick's Theorem]
    \label{thm:Picks_Theorem}
    \uses{def:Simple_Polygon}
    \uses{lem:Additivity_of_Area}
    \uses{lem:Area_of_a_General_Triangle}
    \uses{lem:Additivity_of_Picks_Formula}
    \uses{lem:Polygon_Triangulation}
    For a simple polygon with integer vertices, the area $A$ is given by the formula $A = i + \frac{b}{2} - 1$, where $i$ is the number of integer points in the interior of the polygon and $b$ is the number of integer points on its boundary.
\end{theorem}

\begin{proof}
    Let $P$ be a simple polygon. By lemma \ref{lem:Polygon_Triangulation}, $P$ can be partitioned into a set of $n-2$ non-overlapping triangles $\{T_j\}$.
    Let $F(S) = i_S + \frac{b_S}{2} - 1$.
    From lemma \ref{lem:Area_of_a_General_Triangle}, the formula holds for each triangle, so $A(T_j) = F(T_j)$ for all $j$.
    By lemma \ref{lem:Additivity_of_Area}, the total area of the polygon is the sum of the areas of these triangles: $A(P) = \sum_{j} A(T_j)$.
    By repeated application of lemma \ref{lem:Additivity_of_Picks_Formula}, the function $F$ is additive over the union of these triangles, so $F(P) = \sum_{j} F(T_j)$.
    Combining these facts, we have:
    \[
        A(P) = \sum_{j} A(T_j) = \sum_{j} F(T_j) = F(P).
    \]
    Therefore, $A(P) = i_P + \frac{b_P}{2} - 1$.
\end{proof}

\chapter{Rouch\'{e}--Capelli Theorem}

\begin{definition}[Coefficient Matrix]
    \label{def:Coefficient_Matrix}
    For a system of linear equations with $n$ variables and $m$ equations,
    $$
    \begin{cases}
        a_{11}x_1 + a_{12}x_2 + \dots + a_{1n}x_n = b_1 \\
        a_{21}x_1 + a_{22}x_2 + \dots + a_{2n}x_n = b_2 \\
        \vdots \\
        a_{m1}x_1 + a_{m2}x_2 + \dots + a_{mn}x_n = b_m
    \end{cases}
    $$
    the coefficient matrix, denoted by $A$, is the $m \times n$ matrix containing the coefficients of the variables:
    $$ A = \begin{bmatrix} a_{11} & a_{12} & \dots & a_{1n} \\ a_{21} & a_{22} & \dots & a_{2n} \\ \vdots & \vdots & \ddots & \vdots \\ a_{m1} & a_{m2} & \dots & a_{mn} \end{bmatrix} $$
\end{definition}

\begin{definition}[Augmented Matrix]
    \label{def:Augmented_Matrix}
    \uses{def:Coefficient_Matrix}
    Given a system of linear equations, its augmented matrix, denoted by $[A|b]$, is the coefficient matrix $A$ appended with a column vector $b$ representing the constant terms:
    $$ [A|b] = \left[\begin{array}{cccc|c} a_{11} & a_{12} & \dots & a_{1n} & b_1 \\ a_{21} & a_{22} & \dots & a_{2n} & b_2 \\ \vdots & \vdots & \ddots & \vdots & \vdots \\ a_{m1} & a_{m2} & \dots & a_{mn} & b_m \end{array}\right] $$
\end{definition}

\begin{definition}[Rank of a Matrix]
    \label{def:Rank_of_a_Matrix}
    The rank of a matrix is the maximum number of linearly independent row vectors in the matrix. It is also the number of non-zero rows in its reduced row echelon form.
\end{definition}

\begin{lemma}[Invariance of Solution Set under Row Operations]
    \label{lem:Invariance_of_Solution_Set}
    Applying elementary row operations to an augmented matrix does not change the set of solutions of the corresponding system of linear equations.
\end{lemma}

\begin{proof}
    Elementary row operations (swapping rows, multiplying a row by a non-zero scalar, adding a multiple of one row to another) correspond to reversible transformations of the equations in the system. Since each operation can be undone, the set of variable values satisfying the system remains unchanged.
\end{proof}

\begin{lemma}[Rank Invariance under Row Operations]
    \label{lem:Rank_Invariance_under_Row_Operations}
    \uses{def:Rank_of_a_Matrix}
    Let $M$ be a matrix and let $M'$ be a matrix obtained from $M$ by a finite sequence of elementary row operations. Then the row space of $M$ is identical to the row space of $M'$, and consequently, $\text{rank}(M) = \text{rank}(M')$.
\end{lemma}

\begin{proof}
    Elementary row operations do not change the span of the row vectors of a matrix. Swapping rows does not change the set of row vectors. Multiplying a row by a non-zero scalar scales one of the basis vectors of the row space, but the span remains the same. Adding a multiple of one row to another creates a new vector that is a linear combination of existing vectors, which does not leave the original span. Since these operations are reversible, the row space is preserved. As rank is the dimension of the row space, the rank remains invariant.
\end{proof}

\begin{lemma}[Consistency Condition for RREF Systems]
    \label{lem:Consistency_Condition_RREF}
    A system of linear equations whose augmented matrix is in reduced row echelon form is consistent if and only if it does not contain a row of the form $[0 \ 0 \ \dots \ 0 \ | \ 1]$.
\end{lemma}

\begin{proof}
    If the augmented matrix contains a row of the form $[0 \ 0 \ \dots \ 0 \ | \ 1]$, this corresponds to the contradictory equation $0=1$, so the system is inconsistent. Conversely, if no such row exists, we can express each pivot variable in terms of the non-pivot (free) variables and constants. Assigning arbitrary values to the free variables provides a valid solution, so the system is consistent.
\end{proof}

\begin{lemma}[System Consistency and RREF Structure]
    \label{lem:System_Consistency_and_RREF}
    \uses{lem:Invariance_of_Solution_Set}
    \uses{lem:Consistency_Condition_RREF}
    A system of linear equations is consistent if and only if the reduced row echelon form of its augmented matrix does not contain a row of the form $[0 \ 0 \ \dots \ 0 \ | \ 1]$.
\end{lemma}

\begin{proof}
    \uses{lem:Invariance_of_Solution_Set}
    \uses{lem:Consistency_Condition_RREF}
    Let $S$ be the original system. By lemma \ref{lem:Invariance_of_Solution_Set}, the solution set of $S$ is identical to that of the system $S'$ corresponding to its reduced row echelon form. Thus, $S$ is consistent if and only if $S'$ is consistent. By lemma \ref{lem:Consistency_Condition_RREF}, $S'$ is consistent if and only if its augmented matrix has no row of the form $[0 \ 0 \ \dots \ 0 \ | \ 1]$.
\end{proof}

\begin{lemma}[Rank Equality and RREF Structure]
    \label{lem:Rank_Equality_and_RREF}
    \uses{def:Augmented_Matrix}
    \uses{def:Rank_of_a_Matrix}
    \uses{lem:Rank_Invariance_under_Row_Operations}
    Let $A$ be a coefficient matrix and $[A|b]$ be the corresponding augmented matrix. Then $\text{rank}(A) = \text{rank}([A|b])$ if and only if the reduced row echelon form of $[A|b]$ contains no row of the form $[0 \ 0 \ \dots \ 0 \ | \ 1]$.
\end{lemma}

\begin{proof}
    \uses{def:Rank_of_a_Matrix}
    \uses{lem:Rank_Invariance_under_Row_Operations}
    Let $[A'|b']$ be the reduced row echelon form of $[A|b]$. By lemma \ref{lem:Rank_Invariance_under_Row_Operations}, $\text{rank}(A) = \text{rank}(A')$ and $\text{rank}([A|b]) = \text{rank}([A'|b'])$. The rank is the number of non-zero rows. A row in $[A'|b']$ can be non-zero while the corresponding row in $A'$ is zero only if it is of the form $[0 \ 0 \ \dots \ 0 \ | \ c]$ with $c \ne 0$. In RREF, this means $c=1$. Thus, $\text{rank}(A') < \text{rank}([A'|b'])$ if and only if such a row exists. Consequently, $\text{rank}(A) = \text{rank}([A|b])$ if and only if no such row exists.
\end{proof}

\begin{theorem}[Rouch\'{e}--Capelli Theorem on Existence of Solutions]
    \label{thm:Rouche-Capelli_Existence}
    \uses{def:Coefficient_Matrix}
    \uses{def:Augmented_Matrix}
    \uses{def:Rank_of_a_Matrix}
    A system of linear equations has at least one solution if and only if the rank of its coefficient matrix $A$ is equal to the rank of its augmented matrix $[A|b]$.
\end{theorem}

\begin{proof}
    \uses{lem:System_Consistency_and_RREF}
    \uses{lem:Rank_Equality_and_RREF}
    A system is consistent (has a solution) if and only if the reduced row echelon form of its augmented matrix has no row of the form $[0 \ 0 \ \dots \ 0 \ | \ 1]$, by lemma \ref{lem:System_Consistency_and_RREF}. By lemma \ref{lem:Rank_Equality_and_RREF}, this condition is equivalent to $\text{rank}(A) = \text{rank}([A|b])$.
\end{proof}

\begin{lemma}[General Solution Structure]
    \label{lem:General_Solution_Structure}
    Let $A\mathbf{x}=\mathbf{b}$ be a consistent system of linear equations. If $\mathbf{x}_p$ is a particular solution, then the complete solution set is the set $S = \{ \mathbf{x}_p + \mathbf{x}_h \mid A\mathbf{x}_h = \mathbf{0} \}$.
\end{lemma}

\begin{proof}
    Let $\mathbf{x}$ be any solution to $A\mathbf{x}=\mathbf{b}$. Let $\mathbf{x}_h = \mathbf{x} - \mathbf{x}_p$. Then $A\mathbf{x}_h = A(\mathbf{x} - \mathbf{x}_p) = A\mathbf{x} - A\mathbf{x}_p = \mathbf{b} - \mathbf{b} = \mathbf{0}$. So $\mathbf{x}_h$ is a solution to the homogeneous system. Thus, any solution $\mathbf{x}$ can be written as $\mathbf{x}_p + \mathbf{x}_h$. Conversely, let $\mathbf{x}_h$ be any solution to $A\mathbf{x}_h = \mathbf{0}$. Then $A(\mathbf{x}_p + \mathbf{x}_h) = A\mathbf{x}_p + A\mathbf{x}_h = \mathbf{b} + \mathbf{0} = \mathbf{b}$, so $\mathbf{x}_p + \mathbf{x}_h$ is a solution.
\end{proof}

\begin{lemma}[Rank-Nullity Theorem]
    \label{lem:Rank_Nullity_Theorem}
    \uses{def:Rank_of_a_Matrix}
    For an $m \times n$ matrix $A$, the dimension of its column space (rank) plus the dimension of its null space (nullity) equals the number of columns $n$. That is, $\text{rank}(A) + \dim(N(A)) = n$.
\end{lemma}

\begin{proof}
    Let $r = \text{rank}(A)$. The rank of $A$ is the dimension of its column space, which is also the dimension of its row space. In the reduced row echelon form of $A$, there are $r$ pivot columns and $n-r$ non-pivot columns. The non-pivot columns correspond to free variables in the solution to $A\mathbf{x}=\mathbf{0}$. The null space $N(A)$ is the solution set of $A\mathbf{x}=\mathbf{0}$. We can construct a basis for $N(A)$ by setting one free variable to 1 and all others to 0, and solving for the pivot variables. This can be done for each of the $n-r$ free variables, yielding $n-r$ linearly independent vectors that span $N(A)$. Thus, $\dim(N(A)) = n-r$. The theorem follows.
\end{proof}

\begin{lemma}[Structure of the Solution Set]
    \label{lem:Structure_of_Solution_Set}
    \uses{lem:General_Solution_Structure}
    \uses{lem:Rank_Nullity_Theorem}
    If a consistent system of linear equations in $n$ variables has a coefficient matrix $A$ with $\text{rank}(A) = r$, then the solution set is an affine subspace of dimension $n-r$.
\end{lemma}

\begin{proof}
    \uses{lem:General_Solution_Structure}
    \uses{lem:Rank_Nullity_Theorem}
    By lemma \ref{lem:General_Solution_Structure}, the solution set is a translation of the null space $N(A)$ by a particular solution vector $\mathbf{x}_p$. This is the definition of an affine subspace. The dimension of this affine subspace is the dimension of the vector space $N(A)$. By the Rank-Nullity Theorem (lemma \ref{lem:Rank_Nullity_Theorem}), $\dim(N(A)) = n - \text{rank}(A) = n-r$.
\end{proof}

\begin{theorem}[Rouch\'{e}--Capelli Theorem on the Number of Solutions]
    \label{thm:Rouche-Capelli_Number_of_Solutions}
    \uses{def:Rank_of_a_Matrix}
    \uses{thm:Rouche-Capelli_Existence}
    \uses{lem:Structure_of_Solution_Set}
    If a system of linear equations is consistent, with $n$ variables and coefficient matrix rank $r = \text{rank}(A)$, then:
    \begin{enumerate}
        \item If $r=n$, the system has a unique solution.
        \item If $r<n$, the system has infinitely many solutions (assuming the underlying field is infinite).
    \end{enumerate}
\end{theorem}

\begin{proof}
    \uses{lem:Structure_of_Solution_Set}
    As the system is consistent, we can apply lemma \ref{lem:Structure_of_Solution_Set}. The solution set is an affine subspace of dimension $n-r$.
    \begin{enumerate}
        \item If $r=n$, the dimension is $n-n=0$. A zero-dimensional affine subspace is a single point, hence a unique solution.
        \item If $r<n$, the dimension $n-r$ is positive. An affine subspace of positive dimension over an infinite field contains infinitely many points, hence infinitely many solutions.
    \end{enumerate}
\end{proof}

\chapter{Skolem-Noether Theorem}

\begin{definition}[Simple Ring]
    \label{def:Simple_Ring}
    A ring $R$ is called simple if it has no two-sided ideals other than the zero ideal and itself.
\end{definition}

\begin{definition}[Central Simple Algebra]
    \label{def:Central_Simple_Algebra}
    \uses{def:Simple_Ring}
    A finite-dimensional associative algebra $A$ over a field $k$ is a central simple algebra (CSA) if it is a simple ring and its center is precisely the field $k$.
\end{definition}

\begin{definition}[Opposite Algebra]
    \label{def:Opposite_Algebra}
    Let $(A, \cdot)$ be an algebra over a field $k$. The opposite algebra, denoted $A^{\text{op}}$, has the same underlying vector space as $A$ but with multiplication $*$ defined by $a * b = b \cdot a$ for all $a, b \in A$.
\end{definition}

\begin{lemma}[Opposite of a CSA is a CSA]
    \label{lem:Opposite_of_CSA_is_CSA}
    \uses{def:Central_Simple_Algebra}
    \uses{def:Opposite_Algebra}
    Let $B$ be a central simple algebra over a field $k$. Then its opposite algebra $B^{\text{op}}$ is also a central simple algebra over $k$.
\end{lemma}

\begin{proof}
    The center of an algebra is defined by the multiplication operation. An element $z$ is in the center of $B^{\text{op}}$ if $z * b = b * z$ for all $b \in B^{\text{op}}$. By definition of the opposite algebra, this is equivalent to $b \cdot z = z \cdot b$ for all $b \in B$. This is the condition for $z$ to be in the center of $B$. Therefore, $Z(B^{\text{op}}) = Z(B) = k$.
    Furthermore, a two-sided ideal in $B^{\text{op}}$ is a subset $I$ such that for all $i \in I$ and $b \in B$, $i*b \in I$ and $b*i \in I$. This translates to $b \cdot i \in I$ and $i \cdot b \in I$, which means $I$ is a two-sided ideal in $B$. Since $B$ is simple, its only ideals are $\{0\}$ and $B$. Thus, the only ideals of $B^{\text{op}}$ are $\{0\}$ and $B^{\text{op}}$.
    Since $B^{\text{op}}$ is simple and its center is $k$, it is a central simple algebra.
\end{proof}

\begin{lemma}[Tensor Product of Simple Algebras]
    \label{lem:Tensor_Product_is_Simple}
    \uses{def:Simple_Ring}
    \uses{def:Central_Simple_Algebra}
    If $A$ is a simple algebra and $B$ is a central simple algebra over a field $k$, then the tensor product algebra $A \otimes_k B$ is a simple $k$-algebra.
\end{lemma}

\begin{proof}
    This is a standard theorem in the theory of algebras. The proof relies on showing that any non-zero ideal in the tensor product must contain an element of the form $1 \otimes b$ with $b \neq 0$, from which it can be shown the ideal is the entire algebra.
\end{proof}

\begin{lemma}[Module Structure over Simple Algebras]
    \label{lem:Module_Structure_over_Simple_Algebras}
    \uses{def:Simple_Ring}
    Let $C$ be a finite-dimensional simple algebra over a field $k$. Then there exists a unique simple $C$-module $S$ (up to isomorphism). Furthermore, every finite-dimensional $C$-module $V$ is completely reducible, that is, it is isomorphic to a direct sum of a finite number of copies of $S$, written as $V \cong S^n$ for some unique integer $n \ge 0$.
\end{lemma}

\begin{proof}
    This is a fundamental result from the Artin-Wedderburn theorem, which states that a finite-dimensional simple algebra $C$ is isomorphic to a matrix algebra $\operatorname{M}_d(D)$ over a division ring $D$. The vector space $D^d$ is a simple $C$-module, and it is the unique simple module up to isomorphism. The complete reducibility is a property of modules over semisimple algebras, and every simple algebra is semisimple.
\end{proof}

\begin{lemma}[Isomorphism of Equidimensional Modules]
    \label{lem:Isomorphism_of_Equidimensional_Modules}
    \uses{lem:Module_Structure_over_Simple_Algebras}
    Let $C$ be a finite-dimensional simple algebra over a field $k$. Any two finite-dimensional $C$-modules with the same dimension over $k$ are isomorphic.
\end{lemma}

\begin{proof}
    Let $V_1$ and $V_2$ be two finite-dimensional $C$-modules such that $\dim_k(V_1) = \dim_k(V_2)$. By lemma \ref{lem:Module_Structure_over_Simple_Algebras}, there exists a unique simple module $S$, and we can write $V_1 \cong S^{n_1}$ and $V_2 \cong S^{n_2}$ for unique integers $n_1, n_2$.
    Taking dimensions, we have $\dim_k(V_1) = n_1 \cdot \dim_k(S)$ and $\dim_k(V_2) = n_2 \cdot \dim_k(S)$. Since $\dim_k(V_1) = \dim_k(V_2)$ and $\dim_k(S) > 0$, we must have $n_1 = n_2$. Therefore, $V_1 \cong S^{n_1} \cong V_2$, and the modules are isomorphic.
\end{proof}

\begin{lemma}[Structure of Right Module Homomorphisms on an Algebra]
    \label{lem:Right_Module_Homomorphisms_on_Algebra}
    Let $B$ be a unital algebra over a field $k$. Any right $B$-module homomorphism $u: B \to B$ is given by left multiplication by an element of $B$.
\end{lemma}

\begin{proof}
    Let $u: B \to B$ be a right $B$-module homomorphism. This means $u(b_1 b_2) = u(b_1) b_2$ for all $b_1, b_2 \in B$. Let $c = u(1_B)$, where $1_B$ is the multiplicative identity in $B$. For any element $b \in B$, we can write $b = 1_B \cdot b$. Applying the homomorphism property, we get $u(b) = u(1_B \cdot b) = u(1_B) b = c b$. Thus, the action of $u$ is equivalent to left multiplication by the element $c = u(1_B)$.
\end{proof}

\begin{lemma}[Invertibility of Left Multiplication Map]
    \label{lem:Invertibility_of_Left_Multiplication_Map}
    Let $B$ be a finite-dimensional unital algebra over a field $k$. An element $c \in B$ is invertible if and only if the left multiplication map $L_c: B \to B$ defined by $L_c(b) = cb$ is a vector space isomorphism.
\end{lemma}

\begin{proof}
    For a finite-dimensional vector space, a linear map is an isomorphism if and only if it is injective. The map $L_c$ is injective if its kernel is trivial, i.e., $\ker(L_c) = \{b \in B \mid cb = 0\} = \{0\}$.
    If $c$ is invertible, then $cb = 0$ implies $c^{-1}(cb) = c^{-1}0$, which gives $b=0$. So $L_c$ is injective and thus an isomorphism.
    Conversely, if $L_c$ is an isomorphism, it is surjective. Thus, there exists an element $d \in B$ such that $L_c(d) = cd = 1_B$. This shows $c$ has a right inverse. In a finite-dimensional algebra, an element with a right inverse is invertible.
\end{proof}

\begin{theorem}[Skolem-Noether]
    \label{thm:Skolem-Noether}
    \uses{def:Simple_Ring}
    \uses{def:Central_Simple_Algebra}
    Let $A$ be a simple algebra and $B$ be a central simple algebra over a field $k$. If $f, g: A \to B$ are two $k$-algebra homomorphisms, then there exists an invertible element $c \in B$ such that $g(a) = c f(a) c^{-1}$ for all $a \in A$.
\end{theorem}

\begin{proof}
    \uses{def:Opposite_Algebra}
    \uses{lem:Opposite_of_CSA_is_CSA}
    \uses{lem:Tensor_Product_is_Simple}
    \uses{lem:Isomorphism_of_Equidimensional_Modules}
    \uses{lem:Right_Module_Homomorphisms_on_Algebra}
    \uses{lem:Invertibility_of_Left_Multiplication_Map}
    By lemma \ref{lem:Opposite_of_CSA_is_CSA}, $B^{\text{op}}$ is a central simple algebra. By lemma \ref{lem:Tensor_Product_is_Simple}, the algebra $C = A \otimes_k B^{\text{op}}$ is simple.
    We define two right $C$-module structures on the vector space $B$, denoted $B_f$ and $B_g$, via the actions:
    $b \cdot_f (a \otimes y) = f(a) b y$ and $b \cdot_g (a \otimes y) = g(a) b y$ for $b \in B, a \in A, y \in B^{\text{op}}$.
    Since $B_f$ and $B_g$ have the same finite dimension over $k$, lemma \ref{lem:Isomorphism_of_Equidimensional_Modules} guarantees the existence of a $C$-module isomorphism $u: B_f \to B_g$.
    The condition that $u$ is a $C$-module isomorphism is $u(b \cdot_f x) = u(b) \cdot_g x$ for all $x \in C$.
    Let $x = 1_A \otimes y$ for $y \in B^{\text{op}}$. Then $u(b y) = u(b) y$ for all $b, y \in B$. This shows $u$ is a right $B$-module homomorphism on $B$. By lemma \ref{lem:Right_Module_Homomorphisms_on_Algebra}, there is a $c \in B$ such that $u(b) = cb$ for all $b \in B$. Since $u$ is an isomorphism, lemma \ref{lem:Invertibility_of_Left_Multiplication_Map} implies that $c$ is an invertible element of $B$.
    Now let $x = a \otimes 1_{B^{\text{op}}}$ for $a \in A$. The isomorphism condition is $u(f(a)b) = g(a)u(b)$. Substituting $u(b) = cb$, we get $c(f(a)b) = g(a)(cb)$. This holds for all $b \in B$. For $b=1_B$, we have $cf(a) = g(a)c$. Since $c$ is invertible, we can write $g(a) = c f(a) c^{-1}$.
\end{proof}

\chapter{Stolz–Cesàro theorem}

\begin{lemma}[Addition of Strict Inequalities]
    \label{lem:Addition_of_Strict_Inequalities_Property}
    Let $A, B, C, D$ be real numbers. If $A < B$ and $C < D$, then $A + C < B + D$.
\end{lemma}

\begin{proof}
    The inequality $A < B$ is equivalent to $B - A > 0$. Similarly, $C < D$ is equivalent to $D - C > 0$.
    The sum of two positive real numbers is positive, so $(B - A) + (D - C) > 0$.
    By rearranging the terms, we get $(B + D) - (A + C) > 0$.
    This is equivalent to the inequality $A + C < B + D$.
\end{proof}

\begin{lemma}[Multiplication of an Inequality by a Positive Number]
    \label{lem:Multiplication_by_Positive}
    Let $x, y, z$ be real numbers. If $x < y$ and $z > 0$, then $xz < yz$.
\end{lemma}

\begin{proof}
    The inequality $x < y$ is equivalent to $y - x > 0$.
    We are given that $z > 0$. The product of two positive real numbers is positive, so $(y - x)z > 0$.
    By the distributive property, this is $yz - xz > 0$.
    This is equivalent to the inequality $xz < yz$.
\end{proof}

\begin{lemma}[Multiplication of an Inequality by a Negative Number]
    \label{lem:Multiplication_by_Negative}
    Let $x, y, z$ be real numbers. If $x < y$ and $z < 0$, then $xz > yz$.
\end{lemma}

\begin{proof}
    The inequality $x < y$ is equivalent to $y - x > 0$.
    We are given that $z < 0$. The product of a positive real number and a negative real number is negative, so $(y - x)z < 0$.
    By the distributive property, this is $yz - xz < 0$.
    This is equivalent to the inequality $yz < xz$, or $xz > yz$.
\end{proof}

\begin{lemma}[Additivity of Strict Inequalities over a Sum]
    \label{lem:Additivity_of_Strict_Inequalities}
    \uses{lem:Addition_of_Strict_Inequalities_Property}
    Let $(x_k)$ and $(y_k)$ be two sequences of real numbers. If $x_k < y_k$ for all integers $k$ such that $M \le k \le N$, then $\sum_{k=M}^{N} x_k < \sum_{k=M}^{N} y_k$.
\end{lemma}

\begin{proof}
    We prove this by induction on the number of terms, $n = N - M + 1$.
    Base case: For $n=1$, we have $N=M$. The statement is $x_M < y_M$, which is true by hypothesis.
    Inductive step: Assume the statement is true for $n$ terms. Consider a sum over $n+1$ terms, from $M$ to $M+n$. We have:
    $$\sum_{k=M}^{M+n} x_k = \left(\sum_{k=M}^{M+n-1} x_k\right) + x_{M+n}$$
    By the inductive hypothesis, $\sum_{k=M}^{M+n-1} x_k < \sum_{k=M}^{M+n-1} y_k$. By the given condition, $x_{M+n} < y_{M+n}$.
    Using lemma \ref{lem:Addition_of_Strict_Inequalities_Property}, we can add these two inequalities:
    $$\left(\sum_{k=M}^{M+n-1} x_k\right) + x_{M+n} < \left(\sum_{k=M}^{M+n-1} y_k\right) + y_{M+n}$$
    This simplifies to $\sum_{k=M}^{M+n} x_k < \sum_{k=M}^{M+n} y_k$.
    The statement holds for $n+1$ terms. By the principle of mathematical induction, the lemma is true.
\end{proof}

\begin{lemma}[Telescoping Sum]
    \label{lem:Telescoping_Sum}
    For any sequence of real numbers $(x_k)$ and integers $n > M$, the following identity holds: $\sum_{k=M}^{n-1} (x_{k+1} - x_k) = x_n - x_M$.
\end{lemma}

\begin{proof}
    We prove this by induction on $n$.
    Base case: Let $n=M+1$. The sum is $\sum_{k=M}^{M} (x_{k+1} - x_k) = x_{M+1} - x_M$. The right hand side is $x_{M+1} - x_M$. The identity holds.
    Inductive step: Assume the identity holds for some $n > M$. We want to show it holds for $n+1$.
    $$\sum_{k=M}^{n} (x_{k+1} - x_k) = \left(\sum_{k=M}^{n-1} (x_{k+1} - x_k)\right) + (x_{n+1} - x_n)$$
    By the inductive hypothesis, the sum on the right is equal to $x_n - x_M$.
    $$= (x_n - x_M) + (x_{n+1} - x_n) = x_{n+1} - x_M$$
    This is the required formula for $n+1$. By the principle of mathematical induction, the lemma is proven.
\end{proof}

\begin{lemma}[Property of Sequences Diverging to Infinity]
    \label{lem:Property_of_Divergent_Sequences}
    If a sequence $(b_n)$ diverges to $+\infty$, then for any real number $M$, there exists an integer $N$ such that $b_n > M$ for all $n > N$. In particular, the sequence is eventually positive.
\end{lemma}

\begin{proof}
    The statement is the definition of divergence to $+\infty$. For any $M \in \mathbb{R}$, the definition of $\lim_{n \to \infty} b_n = +\infty$ states that there exists an integer $N$ such that for all $n > N$, we have $b_n > M$. To prove that the sequence is eventually positive, we can choose $M=0$. Then there exists an $N$ such that for all $n > N$, $b_n > 0$.
\end{proof}

\begin{lemma}[Property of Monotone Sequences Converging to Zero]
    \label{lem:Property_of_Monotone_Sequences_Converging_to_Zero}
    If a sequence $(b_n)$ is strictly decreasing and $\lim_{n\to\infty} b_n = 0$, then $b_n > 0$ for all $n$.
\end{lemma}

\begin{proof}
    Assume for contradiction that there exists an integer $N$ such that $b_N \leq 0$. Since the sequence $(b_n)$ is strictly decreasing, for all $n > N$, we have $b_n < b_N \leq 0$. This implies that the limit of the sequence, if it exists, must be less than or equal to $b_N$, so $\lim_{n\to\infty} b_n \leq b_N \leq 0$. However, we are given that the limit is 0. This requires $b_N = 0$. But if $b_N = 0$, then $b_{N+1} < b_N = 0$, which means $\lim_{n\to\infty} b_n \leq b_{N+1} < 0$. This contradicts the fact that the limit is 0. Therefore, our initial assumption must be false, and $b_n > 0$ for all $n$.
\end{proof}

\begin{lemma}[Preservation of Non-Strict Inequalities under Limits]
    \label{lem:Limit_Inequality_Property}
    Let $(x_n)$ and $(y_n)$ be convergent sequences. If there exists an integer $N$ such that $x_n \leq y_n$ for all $n > N$, then $\lim_{n \to \infty} x_n \leq \lim_{n \to \infty} y_n$.
\end{lemma}

\begin{proof}
    Let $L_x = \lim_{n \to \infty} x_n$ and $L_y = \lim_{n \to \infty} y_n$. Assume for contradiction that $L_x > L_y$. Let $\epsilon = \frac{L_x - L_y}{2}$. Since $\epsilon > 0$, by the definition of a limit, there exist integers $N_x$ and $N_y$ such that:
    For $n > N_x$, $|x_n - L_x| < \epsilon \implies x_n > L_x - \epsilon$.
    For $n > N_y$, $|y_n - L_y| < \epsilon \implies y_n < L_y + \epsilon$.
    For any $n > \max(N, N_x, N_y)$, we have both conditions satisfied. Substituting the definition of $\epsilon$:
    $x_n > L_x - \frac{L_x - L_y}{2} = \frac{L_x + L_y}{2}$.
    $y_n < L_y + \frac{L_x - L_y}{2} = \frac{L_x + L_y}{2}$.
    This implies $y_n < \frac{L_x + L_y}{2} < x_n$, so $y_n < x_n$. This contradicts the hypothesis that $x_n \leq y_n$ for all $n>N$. Thus, our assumption must be false, and we must have $L_x \leq L_y$.
\end{proof>

\begin{lemma}[Bounding the sequence difference]
    \label{lem:Bounding_the_sequence_difference}
    \uses{lem:Multiplication_by_Positive}
    \uses{lem:Additivity_of_Strict_Inequalities}
    \uses{lem:Telescoping_Sum}
    Let $(a_n)_{n\geq 1}$ and $(b_n)_{n\geq 1}$ be two sequences of real numbers. Assume that $(b_n)_{n\geq 1}$ is a strictly increasing sequence such that $b_n \to +\infty$.
    If $\lim_{n\to\infty} \frac{a_{n+1}-a_n}{b_{n+1}-b_n} = l$, where $l$ is a finite real number, then for any $\epsilon > 0$, there exists an integer $\nu$ such that for all $n > \nu+1$, the following inequality holds:
    $$(l - \epsilon/2)(b_n - b_{\nu+1}) + a_{\nu+1} < a_n < (l + \epsilon/2)(b_n - b_{\nu+1}) + a_{\nu+1}$$
\end{lemma}

\begin{proof}
    Given $\lim_{n\to\infty} \frac{a_{n+1}-a_n}{b_{n+1}-b_n} = l$, for any $\epsilon/2 > 0$, there exists an integer $\nu$ such that for all $k > \nu$:
    $$l - \frac{\epsilon}{2} < \frac{a_{k+1}-a_k}{b_{k+1}-b_k} < l + \frac{\epsilon}{2}$$
    Since $(b_n)$ is strictly increasing, $b_{k+1} - b_k > 0$. Using lemma \ref{lem:Multiplication_by_Positive}, we can multiply by the positive term $b_{k+1}-b_k$ without changing the direction of the inequalities:
    $$(l - \frac{\epsilon}{2})(b_{k+1}-b_k) < a_{k+1}-a_k < (l + \frac{\epsilon}{2})(b_{k+1}-b_k)$$
    This holds for all integers $k$ from $\nu+1$ to $n-1$. Applying lemma \ref{lem:Additivity_of_Strict_Inequalities}, we sum the inequalities:
    $$\sum_{k=\nu+1}^{n-1} (l - \frac{\epsilon}{2})(b_{k+1}-b_k) < \sum_{k=\nu+1}^{n-1} (a_{k+1}-a_k) < \sum_{k=\nu+1}^{n-1} (l + \frac{\epsilon}{2})(b_{k+1}-b_k)$$
    Using lemma \ref{lem:Telescoping_Sum} on the sums involving $a_k$ and $b_k$, we get:
    $$(l - \frac{\epsilon}{2})(b_n - b_{\nu+1}) < a_n - a_{\nu+1} < (l + \frac{\epsilon}{2})(b_n - b_{\nu+1})$$
    Adding $a_{\nu+1}$ to all parts of the inequality gives the desired result.
\end{proof}

\begin{lemma}[Vanishing of the residual term]
    \label{lem:Vanishing_of_the_residual_term}
    \uses{lem:Property_of_Divergent_Sequences}
    Let $(b_n)_{n\geq 1}$ be a sequence such that $b_n \to +\infty$. For any constant $C$, we have:
    $$\lim_{n\to\infty} \frac{C}{b_n} = 0$$
\end{lemma}
\begin{proof}
    Given any $\epsilon > 0$, we need to find an integer $N$ such that for all $n > N$, $\left|\frac{C}{b_n} - 0\right| < \epsilon$.
    This inequality is equivalent to $\frac{|C|}{|b_n|} < \epsilon$, which simplifies to $|b_n| > \frac{|C|}{\epsilon}$.
    Since $b_n \to +\infty$, by lemma \ref{lem:Property_of_Divergent_Sequences}, for any real number $M$, there exists an $N_1$ such that $b_n > M$ for all $n > N_1$. Let's choose $M = \frac{|C|}{\epsilon}$.
    Also by lemma \ref{lem:Property_of_Divergent_Sequences}, there exists $N_2$ such that $b_n > 0$ for $n > N_2$.
    Let $N = \max(N_1, N_2)$. For $n > N$, we have $b_n > \frac{|C|}{\epsilon}$ and $b_n>0$. Thus $|b_n| = b_n > \frac{|C|}{\epsilon}$.
    This satisfies the condition for the limit to be 0.
\end{proof}

\begin{theorem}[Stolz–Cesàro Theorem for the */∞ case]
    \label{thm:Stolz-Cesaro_Theorem_for_the_infinity_case}
    \uses{lem:Bounding_the_sequence_difference}
    \uses{lem:Vanishing_of_the_residual_term}
    \uses{lem:Property_of_Divergent_Sequences}
    Let $(a_n)_{n\geq 1}$ and $(b_n)_{n\geq 1}$ be two sequences of real numbers. Assume that $(b_n)_{n\geq 1}$ is a strictly monotone and divergent sequence. If the following limit exists:
    $$\lim_{n\to\infty} \frac{a_{n+1}-a_n}{b_{n+1}-b_n} = l$$
    Then, the limit $\lim_{n\to\infty} \frac{a_n}{b_n}$ also exists and is equal to $l$.
\end{theorem}

\begin{proof}
    We prove the case where $(b_n)$ is strictly increasing and $b_n \to +\infty$. The case where it is strictly decreasing and tends to $-\infty$ is similar.
    From lemma \ref{lem:Bounding_the_sequence_difference}, for any $\epsilon > 0$, there exists $\nu$ such that for $n > \nu+1$:
    $$(l - \epsilon/2)(b_n - b_{\nu+1}) + a_{\nu+1} < a_n < (l + \epsilon/2)(b_n - b_{\nu+1}) + a_{\nu+1}$$
    Since $b_n \to +\infty$, by lemma \ref{lem:Property_of_Divergent_Sequences}, there exists an $n_0$ such that $b_n > 0$ for all $n > n_0$. For $n > \max(\nu+1, n_0)$, we can divide the inequality by $b_n$:
    $$(l - \epsilon/2)\left(1 - \frac{b_{\nu+1}}{b_n}\right) + \frac{a_{\nu+1}}{b_n} < \frac{a_n}{b_n} < (l + \epsilon/2)\left(1 - \frac{b_{\nu+1}}{b_n}\right) + \frac{a_{\nu+1}}{b_n}$$
    This can be rewritten as:
    $$(l - \epsilon/2) + \frac{a_{\nu+1} - b_{\nu+1}(l - \epsilon/2)}{b_n} < \frac{a_n}{b_n} < (l + \epsilon/2) + \frac{a_{\nu+1} - b_{\nu+1}(l + \epsilon/2)}{b_n}$$
    The terms $\frac{a_{\nu+1} - b_{\nu+1}(l \pm \epsilon/2)}{b_n}$ are of the form $C/b_n$. By lemma \ref{lem:Vanishing_of_the_residual_term}, these terms go to 0 as $n \to \infty$.
    Thus, there exist $N_1$ and $N_2$ such that for $n>N_1$, $\left|\frac{a_{\nu+1} - b_{\nu+1}(l - \epsilon/2)}{b_n}\right| < \epsilon/2$, and for $n>N_2$, $\left|\frac{a_{\nu+1} - b_{\nu+1}(l + \epsilon/2)}{b_n}\right| < \epsilon/2$.
    For $n > \max(\nu+1, n_0, N_1, N_2)$, we have:
    $$l - \epsilon = (l - \epsilon/2) - \epsilon/2 < \frac{a_n}{b_n} < (l + \epsilon/2) + \epsilon/2 = l + \epsilon$$
    This shows that for any $\epsilon > 0$, there is an $N$ such that for $n>N$, $|\frac{a_n}{b_n} - l| < \epsilon$. Therefore, $\lim_{n\to\infty} \frac{a_n}{b_n} = l$.
\end{proof}

\begin{theorem}[Stolz–Cesàro Theorem for the 0/0 case]
    \label{thm:Stolz-Cesaro_Theorem_for_the_0_case}
    \uses{lem:Multiplication_by_Negative}
    \uses{lem:Additivity_of_Strict_Inequalities}
    \uses{lem:Telescoping_Sum}
    \uses{lem:Limit_Inequality_Property}
    \uses{lem:Property_of_Monotone_Sequences_Converging_to_Zero}
    Let $(a_n)_{n\geq 1}$ and $(b_n)_{n\geq 1}$ be two sequences of real numbers such that $\lim_{n\to\infty} a_n = 0$ and $\lim_{n\to\infty} b_n = 0$. Assume $(b_n)_{n\geq 1}$ is strictly decreasing. If
    $$\lim_{n\to\infty} \frac{a_{n+1}-a_n}{b_{n+1}-b_n} = l$$
    then
    $$\lim_{n\to\infty} \frac{a_n}{b_n} = l$$
\end{theorem}
\begin{proof}
    For any $\epsilon > 0$, there is an integer $N$ such that for all $k \geq N$,
    $$l - \epsilon < \frac{a_{k+1}-a_k}{b_{k+1}-b_k} < l + \epsilon$$
    Since $(b_n)$ is strictly decreasing, $b_{k+1}-b_k < 0$. Using lemma \ref{lem:Multiplication_by_Negative}, we multiply by the negative term $b_{k+1}-b_k$, which reverses the inequalities:
    $$(l+\epsilon)(b_{k+1}-b_k) < a_{k+1}-a_k < (l-\epsilon)(b_{k+1}-b_k)$$
    Let $n > N$. For any $m > n$, we sum from $k=n$ to $m-1$. By lemma \ref{lem:Additivity_of_Strict_Inequalities}:
    $$\sum_{k=n}^{m-1} (l+\epsilon)(b_{k+1}-b_k) < \sum_{k=n}^{m-1} (a_{k+1}-a_k) < \sum_{k=n}^{m-1} (l-\epsilon)(b_{k+1}-b_k)$$
    Using lemma \ref{lem:Telescoping_Sum}, we simplify the sums:
    $$(l+\epsilon)(b_m-b_n) < a_m-a_n < (l-\epsilon)(b_m-b_n)$$
    Now we take the limit as $m \to \infty$. Since $\lim_{m\to\infty} a_m = 0$ and $\lim_{m\to\infty} b_m = 0$, by lemma \ref{lem:Limit_Inequality_Property}, we can take the limit inside the inequalities:
    $$(l+\epsilon)(-b_n) \leq -a_n \leq (l-\epsilon)(-b_n)$$
    By lemma \ref{lem:Property_of_Monotone_Sequences_Converging_to_Zero}, since $(b_n)$ is strictly decreasing to 0, we have $b_n > 0$ for all $n$. Thus, $-b_n < 0$. Dividing the entire inequality by $-b_n$ reverses the inequalities again:
    $$l-\epsilon \leq \frac{a_n}{b_n} \leq l+\epsilon$$
    This can be written as $|\frac{a_n}{b_n} - l| \leq \epsilon$. Since this holds for any $n>N$ and for any $\epsilon > 0$, we conclude that $\lim_{n\to\infty} \frac{a_n}{b_n} = l$.
\end{proof}

\chapter{Thales's Theorem}

\begin{definition}[Angle Inscribed in a Semicircle]
    \label{def:Angle_Inscribed_in_a_Semicircle}
    An angle $\angle ACB$ is inscribed in a semicircle if A, B, and C are distinct points on a circle where the line segment AB is a diameter of the circle.
\end{definition}

\begin{lemma}[Side-Side-Side (SSS) Congruence]
    \label{lem:SSS_Congruence}
    If the three sides of one triangle are equal in length to the corresponding three sides of another triangle, then the two triangles are congruent.
\end{lemma}

\begin{proof}
    This is a fundamental postulate of Euclidean geometry for determining the congruence of triangles.
\end{proof}

\begin{lemma}[Corresponding Parts of Congruent Triangles are Congruent (CPCTC)]
    \label{lem:CPCTC}
    If two triangles are congruent, then their corresponding angles and sides are equal in measure.
\end{lemma}

\begin{proof}
    This follows from the definition of triangle congruence. Congruent triangles can be perfectly superimposed, meaning all their corresponding parts (angles and sides) must be identical.
\end{proof}

\begin{lemma}[Alternate Interior Angles Theorem]
    \label{lem:Alternate_Interior_Angles}
    If two parallel lines are intersected by a transversal, then the alternate interior angles are equal in measure.
\end{lemma}

\begin{proof}
    This is a theorem in Euclidean geometry that follows from the parallel postulate.
\end{proof}

\begin{lemma}[Angles on a Straight Line]
    \label{lem:Angles_on_a_Straight_Line}
    The sum of the measures of adjacent angles that form a straight line is $180^\circ$.
\end{lemma}

\begin{proof}
    This is a fundamental postulate of geometry, defining a straight angle as having a measure of $180^\circ$.
\end{proof}

\begin{lemma}[Angle Addition Postulate]
    \label{lem:Angle_Addition_Postulate}
    If a point O lies in the interior of an angle $\angle ACB$, then the measure of the angle is the sum of its non-overlapping component parts: $\angle ACB = \angle ACO + \angle OCB$.
\end{lemma}

\begin{proof}
    This postulate is a fundamental property of angle measurement in Euclidean geometry. It defines how angles are composed from smaller, adjacent angles sharing a common vertex and a common ray between them.
\end{proof}

\begin{lemma}[Isosceles Triangle Theorem]
    \label{lem:Isosceles_Triangle_Theorem}
    \uses{lem:SSS_Congruence}
    \uses{lem:CPCTC}
    If two sides of a triangle are equal in length, then the angles opposite those sides are equal in measure.
\end{lemma}

\begin{proof}
    Let $\triangle ABC$ be a triangle with sides $AC = BC$. Let $M$ be the midpoint of the segment $AB$. Then the segment $CM$ is a median. Consider the two triangles $\triangle AMC$ and $\triangle BMC$. We have $AC = BC$ (given), $AM = BM$ (by construction of midpoint), and $CM = CM$ (common side). By the Side-Side-Side congruence criterion stated in lemma \ref{lem:SSS_Congruence}, $\triangle AMC \cong \triangle BMC$. Since the triangles are congruent, their corresponding angles are equal, as per lemma \ref{lem:CPCTC}. In particular, $\angle CAM$ corresponds to $\angle CBM$. Therefore, $\angle CAB = \angle CBA$.
\end{proof}

\begin{lemma}[Sum of Angles in a Triangle]
    \label{lem:Sum_of_Angles_in_a_Triangle}
    \uses{lem:Alternate_Interior_Angles}
    \uses{lem:Angles_on_a_Straight_Line}
    The sum of the measures of the interior angles of any triangle in Euclidean geometry is $180^\circ$.
\end{lemma}

\begin{proof}
    Given $\triangle ABC$, construct a line $L$ passing through vertex C and parallel to the side $AB$. Let angles $\alpha$ and $\beta$ be the angles that line $L$ makes with segments $AC$ and $BC$, respectively, on the same side of $L$ as the triangle. By the alternate interior angles theorem, lemma \ref{lem:Alternate_Interior_Angles}, we have $\angle CAB = \alpha$ and $\angle CBA = \beta$. The angles $\alpha$, $\beta$, and $\angle ACB$ at vertex C form a straight line. According to lemma \ref{lem:Angles_on_a_Straight_Line}, their sum is $180^\circ$. Therefore, $\alpha + \beta + \angle ACB = 180^\circ$. Substituting back, we get $\angle CAB + \angle CBA + \angle ACB = 180^\circ$.
\end{proof}

\begin{lemma}[Radii Form Isosceles Triangles]
    \label{lem:Radii_Form_Isosceles_Triangles}
    \uses{def:Angle_Inscribed_in_a_Semicircle}
    Let $\angle ACB$ be an angle inscribed in a semicircle with center O and diameter AB. Then the triangles $\triangle OAC$ and $\triangle OBC$ are isosceles triangles.
\end{lemma}

\begin{proof}
    The points A, B, and C lie on the circle, and O is the center. The segments OA, OB, and OC are radii of the circle. By definition, all radii of a circle are equal in length, so $OA = OC$ and $OB = OC$. A triangle with at least two equal sides is an isosceles triangle. Therefore, $\triangle OAC$ (with sides $OA=OC$) and $\triangle OBC$ (with sides $OB=OC$) are isosceles.
\end{proof}

\begin{lemma}[Equality of Base Angles in Constructed Triangles]
    \label{lem:Equality_of_Base_Angles}
    \uses{lem:Isosceles_Triangle_Theorem}
    \uses{lem:Radii_Form_Isosceles_Triangles}
    If $\angle ACB$ is an angle inscribed in a semicircle with center O, then the base angles of the constructed triangles $\triangle OAC$ and $\triangle OBC$ are equal. Specifically, $\angle OAC = \angle OCA$ and $\angle OBC = \angle OCB$.
\end{lemma}

\begin{proof}
    From lemma \ref{lem:Radii_Form_Isosceles_Triangles}, we know that $\triangle OAC$ and $\triangle OBC$ are isosceles triangles. According to lemma \ref{lem:Isosceles_Triangle_Theorem}, the angles opposite the equal sides are equal. In $\triangle OAC$, since $OA=OC$, we have $\angle OCA = \angle OAC$. In $\triangle OBC$, since $OB=OC$, we have $\angle OCB = \angle OBC$.
\end{proof}

\begin{lemma}[Sum of Angles as a Multiple of the Inscribed Angle]
    \label{lem:Sum_as_Multiple_of_Inscribed_Angle}
    \uses{lem:Angle_Addition_Postulate}
    \uses{lem:Equality_of_Base_Angles}
    If $\angle ACB$ is an angle inscribed in a semicircle, then the sum of the interior angles of $\triangle ABC$ is equal to $2 \times \angle ACB$.
\end{lemma}

\begin{proof}
    The sum of the interior angles of $\triangle ABC$ is $\angle CAB + \angle ABC + \angle BCA$.
    From lemma \ref{lem:Equality_of_Base_Angles}, we have $\angle CAB = \angle OCA$ and $\angle ABC = \angle OCB$.
    Substituting these into the sum gives: $\angle OCA + \angle OCB + \angle BCA$.
    According to lemma \ref{lem:Angle_Addition_Postulate}, $\angle BCA = \angle OCA + \angle OCB$.
    Substituting this gives $\angle BCA + \angle BCA = 2 \times \angle BCA$, also denoted as $2 \times \angle ACB$.
\end{proof}

\begin{theorem}[Thales's Theorem]
    \label{thm:Thales_Theorem}
    \uses{def:Angle_Inscribed_in_a_Semicircle}
    \uses{lem:Sum_of_Angles_in_a_Triangle}
    \uses{lem:Sum_as_Multiple_of_Inscribed_Angle}
    If A, B, and C are distinct points on a circle where the line segment AB is a diameter, then the angle $\angle ACB$ is a right angle.
\end{theorem}

\begin{proof}
    According to lemma \ref{lem:Sum_of_Angles_in_a_Triangle}, the sum of the interior angles of $\triangle ABC$ is $180^\circ$.
    From lemma \ref{lem:Sum_as_Multiple_of_Inscribed_Angle}, this same sum is equal to $2 \times \angle ACB$.
    By equating these two expressions, we obtain the equation $2 \times \angle ACB = 180^\circ$.
    Dividing both sides of the equation by 2 gives $\angle ACB = 90^\circ$.
    Therefore, $\angle ACB$ is a right angle.
\end{proof}
